% ============================================================
% HSBI --- Quantitative Research Methods
% Mock Examination 2 (Lectures 5--8, ISLP Chapters 4--6)
% Prof. Dr. Christoph Weisser --- Summer Semester 2026
% ------------------------------------------------------------
% SINGLE SOURCE with solutions toggle:
%   Exam only:
%     pdflatex -jobname=Mock_Exam_2 mock_exam_2.tex
%   With solutions:
%     pdflatex -jobname=Mock_Exam_2_Solutions "\def\withsolutions{1}% ============================================================
% HSBI --- Quantitative Research Methods
% Mock Examination 2 (Lectures 5--8, ISLP Chapters 4--6)
% Prof. Dr. Christoph Weisser --- Summer Semester 2026
% ------------------------------------------------------------
% SINGLE SOURCE with solutions toggle:
%   Exam only:
%     pdflatex -jobname=Mock_Exam_2 mock_exam_2.tex
%   With solutions:
%     pdflatex -jobname=Mock_Exam_2_Solutions "\def\withsolutions{1}% ============================================================
% HSBI --- Quantitative Research Methods
% Mock Examination 2 (Lectures 5--8, ISLP Chapters 4--6)
% Prof. Dr. Christoph Weisser --- Summer Semester 2026
% ------------------------------------------------------------
% SINGLE SOURCE with solutions toggle:
%   Exam only:
%     pdflatex -jobname=Mock_Exam_2 mock_exam_2.tex
%   With solutions:
%     pdflatex -jobname=Mock_Exam_2_Solutions "\def\withsolutions{1}% ============================================================
% HSBI --- Quantitative Research Methods
% Mock Examination 2 (Lectures 5--8, ISLP Chapters 4--6)
% Prof. Dr. Christoph Weisser --- Summer Semester 2026
% ------------------------------------------------------------
% SINGLE SOURCE with solutions toggle:
%   Exam only:
%     pdflatex -jobname=Mock_Exam_2 mock_exam_2.tex
%   With solutions:
%     pdflatex -jobname=Mock_Exam_2_Solutions "\def\withsolutions{1}\input{mock_exam_2.tex}"
% ============================================================
\documentclass[11pt,a4paper]{article}
\usepackage[utf8]{inputenc}
\usepackage[T1]{fontenc}
\usepackage[english]{babel}
\usepackage[margin=2.2cm]{geometry}
\usepackage{amsmath,amssymb}
\usepackage{booktabs,array}
\usepackage{enumitem}
\usepackage{xcolor}
\usepackage{tcolorbox}
\tcbuselibrary{skins,breakable}

\setlength{\parindent}{0pt}

% Teal solution box (only shown when \withsolutions is defined)
\newtcolorbox{solutionbox}{enhanced,breakable,colback=teal!5,colframe=teal!55!black,
  boxrule=0pt,leftrule=2.5pt,arc=2pt,fonttitle=\bfseries,title=Solution,
  left=2mm,right=2mm,top=1mm,bottom=1.5mm}

\newcommand{\problem}[2]{\par\vspace{7mm}\noindent{\large\textbf{Problem #1}}%
  \hfill{\large\textbf{(#2 points)}}\par\nopagebreak\vspace{0.6mm}%
  \noindent\rule{\textwidth}{0.6pt}\par\nopagebreak\vspace{2.5mm}}
\newcommand{\pp}[1]{\textbf{[#1~P]}\enspace}

\begin{document}

% ------------------------------------------------------------
% Header block
% ------------------------------------------------------------
\begin{center}
  {\Large\textbf{HSBI --- Bielefeld University of Applied Sciences and Arts}}\\[1.2mm]
  {\large Quantitative Research Methods}\\[1mm]
  {\normalsize Summer Semester 2026 \quad$\cdot$\quad Examiner: Prof.\ Dr.\ Christoph Weisser}\\[4.5mm]
  {\LARGE\textbf{Mock Examination 2 (Lectures 5--8)}}\\[1.2mm]
  {\normalsize ISLP Chapters 4--6: Classification $\cdot$ Resampling Methods $\cdot$ Model Selection and Regularisation}%
  \ifdefined\withsolutions\\[2.2mm]{\large\textcolor{teal!55!black}{\textbf{--- Version with detailed solutions ---}}}\fi
\end{center}

\vspace{2.5mm}
\begin{center}
\begin{tabular}{@{}ll@{}}
  \textbf{Duration:} & \textbf{90 minutes}\\
  \textbf{Total credit:} & \textbf{90 points}\\
  \textbf{Allowed aids:} & non-programmable calculator; one handwritten A4 sheet of notes (both sides)\\
\end{tabular}
\end{center}

\vspace{3mm}
Name: \rule[-1mm]{0.38\textwidth}{0.4pt}\hfill Matriculation no.: \rule[-1mm]{0.24\textwidth}{0.4pt}

\vspace{4.5mm}
\begin{center}
\renewcommand{\arraystretch}{1.35}
\begin{tabular}{|l|c|c|c|c|c|c|c|}
\hline
\textbf{Problem} & \textbf{P1} & \textbf{P2} & \textbf{P3} & \textbf{P4} & \textbf{P5} & \textbf{P6} & \textbf{Total}\\\hline
Maximum points & 18 & 12 & 15 & 18 & 12 & 15 & \textbf{90}\\\hline
Achieved points & \hspace{9mm} & \hspace{9mm} & \hspace{9mm} & \hspace{9mm} & \hspace{9mm} & \hspace{9mm} & \hspace{11mm}\\\hline
\end{tabular}
\end{center}

\vspace{2.5mm}
\textbf{Instructions}
\begin{itemize}[nosep,leftmargin=5.5mm]
  \item Justify all answers. Numerical results without a derivation or a brief reasoning receive no credit.
  \item Round final numerical results to 3 decimal places unless stated otherwise. Small deviations caused by carrying rounded intermediate results are accepted.
  \item The number of points per problem roughly equals the recommended working time in minutes.
\end{itemize}
\vspace{1mm}\noindent\rule{\textwidth}{0.6pt}

% ============================================================
\problem{1: Logistic regression by hand}{18}

A bank models the probability that a credit-card customer \emph{defaults} ($Y=1$) as a
function of the monthly card balance $x$ (in EUR; balances in this portfolio range from
0 to 400 EUR). A logistic regression fitted by maximum likelihood gave
\[
\hat{p}(x)=\Pr(Y=1\mid X=x)
  =\frac{e^{\hat{\beta}_0+\hat{\beta}_1 x}}{1+e^{\hat{\beta}_0+\hat{\beta}_1 x}},
\qquad
\hat{\beta}_0=-4,\quad \hat{\beta}_1=0.02 .
\]

\begin{enumerate}[label=(\alph*),leftmargin=7mm,itemsep=2.5mm]
  \item \pp{4} Compute the predicted default probability $\hat{p}(x)$ for a customer with
    balance $x=100$ EUR and for a customer with balance $x=250$ EUR.
  \item \pp{4} For the same two customers, compute the \emph{odds} of default
    $\hat{p}(x)/(1-\hat{p}(x))$ and the \emph{log-odds} (logit). Interpret the odds of the
    customer with $x=250$ in one sentence.
  \item \pp{3} Determine the balance $x^{*}$ at which the predicted default probability
    equals exactly $0.5$. Show your derivation.
  \item \pp{4} Interpret the estimated slope on the odds scale: by which \emph{factor} do the
    estimated odds of default change (i) per additional EUR of balance and (ii) per
    additional 100 EUR of balance? Why is the statement ``$\hat{\beta}_1=0.02$ means the
    default \emph{probability} rises by 0.02 per EUR'' wrong?
  \item \pp{3} Give two reasons why an ordinary linear regression of the 0/1 response $Y$
    on $x$ would be inappropriate here, even though it can be computed mechanically.
\end{enumerate}

\ifdefined\withsolutions
\begin{solutionbox}
\textbf{(a) Predicted probabilities.} The linear predictor (log-odds) is
$\hat{\eta}(x)=\hat{\beta}_0+\hat{\beta}_1x=-4+0.02x$.
\[
x=100:\ \hat{\eta}=-4+2=-2,\qquad
\hat{p}(100)=\frac{e^{-2}}{1+e^{-2}}=\frac{0.1353}{1.1353}=0.119 .
\]
\[
x=250:\ \hat{\eta}=-4+5=1,\qquad
\hat{p}(250)=\frac{e^{1}}{1+e^{1}}=\frac{2.7183}{3.7183}=0.731 .
\]

\medskip
\textbf{(b) Odds and log-odds.} Since
$\hat{p}/(1-\hat{p})=e^{\hat{\eta}}$ and $\log\big(\hat{p}/(1-\hat{p})\big)=\hat{\eta}$:
\begin{center}
\begin{tabular}{lcc}
\toprule
 & $x=100$ & $x=250$\\
\midrule
log-odds $\hat{\eta}(x)$ & $-2.000$ & $1.000$\\
odds $e^{\hat{\eta}(x)}$ & $e^{-2}=0.135$ & $e^{1}=2.718$\\
\bottomrule
\end{tabular}
\end{center}
Check via (a): $0.119/0.881=0.135$ and $0.731/0.269=2.718$. \checkmark\
\emph{Interpretation for $x=250$:} the odds are about $2.72$ to 1, i.e.\ default is
estimated to be about $2.7$ times as probable as non-default for this customer.

\medskip
\textbf{(c) Balance with $\hat{p}=0.5$.}
$\hat{p}(x^{*})=0.5$ $\iff$ odds $=1$ $\iff$ log-odds $=0$:
\[
-4+0.02\,x^{*}=0
\quad\Longrightarrow\quad
x^{*}=\frac{4}{0.02}=200\ \text{EUR}.
\]
Below 200 EUR the model predicts ``no default'' as more likely, above 200 EUR ``default''.

\medskip
\textbf{(d) Odds-scale interpretation.}
Per additional EUR the odds are multiplied by
$e^{\hat{\beta}_1}=e^{0.02}=1.020$, i.e.\ they increase by about $2.0\,\%$.
Per additional 100 EUR they are multiplied by
$e^{100\hat{\beta}_1}=e^{2}=7.389$ (an increase of about $639\,\%$).
The quoted statement is wrong because $\hat{\beta}_1$ is the change in the
\emph{log-odds} per EUR, not in the probability: the probability change
$\partial\hat{p}/\partial x=\hat{\beta}_1\hat{p}(x)\bigl(1-\hat{p}(x)\bigr)$ depends on $x$
--- the S-shaped curve is steepest near $\hat{p}=0.5$ and nearly flat in the tails, so there
is no constant ``probability effect per EUR''.

\medskip
\textbf{(e) Why not linear regression?}
(1) A straight line $\hat{\beta}_0+\hat{\beta}_1x$ is unbounded, so for small or large
balances it produces ``probabilities'' below 0 or above 1, which are not sensible
estimates of $\Pr(Y=1\mid x)$; the logistic function is guaranteed to stay in $[0,1]$.
(2) The 0/1 response violates the linear-model error assumptions: the errors are not
normal and their variance $p(x)(1-p(x))$ depends on $x$ (heteroscedasticity), so the usual
inference is unreliable. (Also, with more than two classes a numeric coding of $Y$ would
impose an artificial ordering and equal spacing between classes.)
\end{solutionbox}
\fi

% ============================================================
\problem{2: Discriminant analysis and naive Bayes --- concepts}{12}

\begin{enumerate}[label=(\alph*),leftmargin=7mm,itemsep=2.5mm]
  \item \pp{4} Linear discriminant analysis (LDA) and quadratic discriminant analysis (QDA)
    both model the class-conditional densities $f_k(x)$ of the predictors as multivariate
    normal. State precisely the assumption about the covariance matrices that distinguishes
    the two methods, and give the number of covariance \emph{matrices} each method estimates
    when there are $K$ classes.
  \item \pp{3} Using the bias--variance trade-off, explain in which situations LDA tends to
    outperform QDA on test data, and in which situations QDA tends to win. Mention the role
    of the sample size $n$ and of the true class covariances.
  \item \pp{3} State the key assumption of the \emph{naive Bayes} classifier about the
    predictors within each class, and explain why this assumption can improve test
    performance when the number of predictors $p$ is large relative to $n$ --- even though
    the assumption is rarely exactly true.
  \item \pp{2} Which of the three classifiers (LDA, QDA, naive Bayes) produce \emph{linear}
    decision boundaries in $x$, and which produce \emph{quadratic} or more general
    non-linear boundaries?
\end{enumerate}

\ifdefined\withsolutions
\begin{solutionbox}
\textbf{(a) Covariance assumptions.}
Both methods assume $X\mid Y=k\sim\mathcal{N}(\mu_k,\Sigma_k)$.
\emph{LDA} additionally assumes a covariance matrix that is \emph{shared} by all classes:
$\Sigma_1=\dots=\Sigma_K=\Sigma$; it therefore estimates \emph{one} common $p\times p$
covariance matrix. \emph{QDA} allows each class its \emph{own} covariance matrix
$\Sigma_k$ and estimates $K$ covariance matrices (one per class), i.e.\ roughly
$K\,p(p+1)/2$ covariance parameters instead of $p(p+1)/2$.

\medskip
\textbf{(b) When does which method win?}
QDA is the more flexible method (many more parameters): \emph{low bias, high variance};
LDA is more restrictive: \emph{higher bias if the $\Sigma_k$ truly differ, but much lower
variance}. Hence LDA tends to win when $n$ is small (or $p$ large), so that reducing
variance is crucial, and when the true class covariances are (approximately) equal ---
then LDA's assumption costs little bias. QDA tends to win when $n$ is large, so its extra
parameters can be estimated reliably, and when the class covariances clearly differ, so
LDA's shared-$\Sigma$ assumption creates substantial bias.

\medskip
\textbf{(c) Naive Bayes.}
Assumption: \emph{within each class the $p$ predictors are independent}, i.e.\ the joint
density factorises as $f_k(x)=f_{k1}(x_1)\cdot f_{k2}(x_2)\cdots f_{kp}(x_p)$.
Instead of a full $p$-dimensional joint density (which would require estimating
associations between all predictors, e.g.\ $p(p+1)/2$ covariance parameters per class),
only $p$ one-dimensional marginal densities per class must be estimated. This drastic
reduction in the number of parameters lowers the \emph{variance} of the classifier.
When $p$ is large relative to $n$, the variance saving usually outweighs the bias
introduced by (falsely) ignoring within-class dependence --- another instance of the
bias--variance trade-off.

\medskip
\textbf{(d) Shape of the decision boundaries.}
\emph{LDA:} linear boundaries (the discriminant functions are linear in $x$; the quadratic
term cancels because $\Sigma$ is shared). \emph{QDA:} quadratic boundaries (each class has
its own quadratic form). \emph{Naive Bayes:} generally non-linear boundaries of flexible
shape --- the log-odds are an \emph{additive} function $\sum_j g_j(x_j)$ of the individual
predictors (no interaction terms); only in special cases (e.g.\ Gaussian marginals with
class-shared variances) does it reduce to a linear rule.
\end{solutionbox}
\fi

% ============================================================
\problem{3: Confusion matrix and ROC curve}{15}

A logistic regression classifier predicts credit-card default on an evaluation set of
$n=1{,}000$ customers, using the default threshold: predict ``default'' if
$\hat{p}(\text{default}\mid x)>0.5$. The result is the following confusion matrix:

\begin{center}
\renewcommand{\arraystretch}{1.15}
\begin{tabular}{llcc|c}
 & & \multicolumn{2}{c|}{\textbf{True status}} & \\
 & & Default & No default & Total\\
\cmidrule{2-5}
\textbf{Predicted} & Default    & 60 & 45  & 105\\
                   & No default & 40 & 855 & 895\\
\cmidrule{2-5}
 & Total & 100 & 900 & 1{,}000\\
\end{tabular}
\end{center}

\begin{enumerate}[label=(\alph*),leftmargin=7mm,itemsep=2.5mm]
  \item \pp{7} Identify TP, FP, FN and TN, and compute (i) the overall accuracy and the
    overall error rate, (ii) the sensitivity (recall, true positive rate), (iii) the
    specificity, (iv) the precision, and (v) the false positive rate. Interpret sensitivity
    and precision in one sentence each in the context of this application.
  \item \pp{4} The bank now lowers the classification threshold from $0.5$ to $0.2$
    (predict ``default'' if $\hat{p}>0.2$). Explain --- without further computation --- how
    sensitivity and specificity will change, and give a business reason why the bank might
    deliberately accept this trade-off.
  \item \pp{4} The ROC curve plots which two quantities against each other, as what is
    varied? What does an area under the curve (AUC) of $0.5$ mean, and what does an AUC of
    $1$ mean?
\end{enumerate}

\ifdefined\withsolutions
\begin{solutionbox}
\textbf{(a) Metrics.} With ``default'' as the positive class:
$\text{TP}=60$, $\text{FP}=45$, $\text{FN}=40$, $\text{TN}=855$.
\begin{align*}
\text{(i) accuracy} &= \frac{\text{TP}+\text{TN}}{n}=\frac{60+855}{1000}=\frac{915}{1000}=0.915,
\qquad \text{error rate}=1-0.915=0.085,\\[0.5mm]
\text{(ii) sensitivity} &= \frac{\text{TP}}{\text{TP}+\text{FN}}=\frac{60}{100}=0.600,\\[0.5mm]
\text{(iii) specificity} &= \frac{\text{TN}}{\text{TN}+\text{FP}}=\frac{855}{900}=0.950,\\[0.5mm]
\text{(iv) precision} &= \frac{\text{TP}}{\text{TP}+\text{FP}}=\frac{60}{105}=0.571,\\[0.5mm]
\text{(v) FPR} &= \frac{\text{FP}}{\text{FP}+\text{TN}}=\frac{45}{900}=0.050=1-\text{specificity}.
\end{align*}
\emph{Sensitivity:} the classifier detects $60.0\,\%$ of the customers who actually
default --- $40\,\%$ of true defaulters are missed.
\emph{Precision:} of the customers flagged as defaulters, only $57.1\,\%$ actually
default; the rest are false alarms.

\medskip
\textbf{(b) Lowering the threshold to 0.2.}
More customers exceed the lower cut-off, so more are predicted ``default'': every true
defaulter with $0.2<\hat{p}\le 0.5$ is now caught, so \emph{sensitivity increases}
(FN $\downarrow$); at the same time more non-defaulters are flagged, so
\emph{specificity decreases} (FP $\uparrow$, FPR $\uparrow$). (The overall error rate will
typically increase, since non-defaulters are the large majority.)
\emph{Business reason:} the costs are asymmetric --- a missed defaulter (FN, a credit
loss) is far more expensive than a false alarm (FP, e.g.\ an unnecessary credit-line
review or customer call). Trading specificity for sensitivity can therefore minimise the
\emph{expected cost} even if the raw error rate rises.

\medskip
\textbf{(c) ROC and AUC.}
The ROC curve plots the \emph{true positive rate} (sensitivity) on the vertical axis
against the \emph{false positive rate} ($1-$specificity) on the horizontal axis, as the
classification \emph{threshold} is varied over all values from 1 down to 0.
$\text{AUC}=0.5$: performance equal to the diagonal, i.e.\ no better than random guessing
--- the score contains no usable ranking information. $\text{AUC}=1$: perfect
classifier --- some threshold separates the two classes without error
(sensitivity $=1$ and specificity $=1$ simultaneously). Equivalently, the AUC is the
probability that a randomly chosen defaulter receives a higher score $\hat{p}$ than a
randomly chosen non-defaulter.
\end{solutionbox}
\fi

% ============================================================
\problem{4: Resampling methods}{18}

\begin{enumerate}[label=(\alph*),leftmargin=7mm,itemsep=2.5mm]
  \item \pp{4} A data set has $n=200$ observations. To estimate the test error of a fixed
    model specification, state how many times the model must be \emph{fitted} for (i) the
    validation-set approach (one 50/50 split), (ii) 10-fold cross-validation, and (iii)
    leave-one-out cross-validation (LOOCV), with a one-sentence justification each.
  \item \pp{4} Compare LOOCV and 10-fold CV as estimators of the true test error with
    respect to \emph{bias} and \emph{variance}. Which method is preferred in practice, and
    why?
  \item \pp{5} A bootstrap sample of size $n$ is drawn with replacement from $n$
    observations. Show that the probability that a fixed observation $i$ is
    \emph{contained} in the bootstrap sample equals $1-(1-1/n)^n$. Evaluate this
    probability for $n=5$, and state the limiting value as $n\to\infty$
    (use $e^{-1}\approx 0.368$).
  \item \pp{5} To quantify the uncertainty of an estimated portfolio weight
    $\hat{\alpha}$, an analyst draws $B=3$ bootstrap samples (in practice $B\approx 1{,}000$;
    $B=3$ is used here only to keep the arithmetic manageable) and obtains
    \[
    \hat{\alpha}^{*1}=0.4,\qquad \hat{\alpha}^{*2}=0.5,\qquad \hat{\alpha}^{*3}=0.6 .
    \]
    Compute the bootstrap estimate of the standard error,
    \[
    \widehat{\operatorname{SE}}_B(\hat{\alpha})
      =\sqrt{\frac{1}{B-1}\sum_{r=1}^{B}\Bigl(\hat{\alpha}^{*r}-\bar{\alpha}^{*}\Bigr)^{2}},
    \qquad
    \bar{\alpha}^{*}=\frac{1}{B}\sum_{r=1}^{B}\hat{\alpha}^{*r},
    \]
    and explain in one sentence what this number estimates. Why is the bootstrap useful
    here although a formula for $\operatorname{SE}(\hat{\alpha})$ may be hard to derive?
\end{enumerate}

\ifdefined\withsolutions
\begin{solutionbox}
\textbf{(a) Number of model fits at $n=200$.}
(i) \emph{Validation set:} \textbf{1} fit --- the model is trained once on the training
half (100 observations) and evaluated on the held-out half.
(ii) \emph{10-fold CV:} \textbf{10} fits --- each of the 10 folds (20 observations) is
held out once, and the model is refitted on the remaining 180 observations each time.
(iii) \emph{LOOCV:} \textbf{200} fits --- each single observation is held out once and
the model is refitted on the other 199. (LOOCV $=$ $n$-fold CV; for least squares a
shortcut formula avoids the refits, but generically $n$ fits are needed.)

\medskip
\textbf{(b) Bias and variance of the CV estimators.}
\emph{Bias:} LOOCV trains on $n-1$ observations, almost the full sample, so it is nearly
unbiased for the true test error; 10-fold CV trains on $0.9\,n$ and therefore slightly
\emph{overestimates} the test error (a bit more bias). \emph{Variance:} LOOCV has the
\emph{higher} variance: the $n$ fitted models are trained on almost identical data sets,
so their errors are strongly positively correlated, and the mean of highly correlated
quantities has high variance; the 10 folds of 10-fold CV overlap less, giving less
correlated fits and a lower-variance estimate. \emph{In practice} $k=5$ or $k=10$ is
preferred: a good bias--variance compromise and much cheaper computationally.

\medskip
\textbf{(c) Bootstrap inclusion probability.}
Each of the $n$ draws is made with replacement, so each draw selects observation $i$ with
probability $1/n$ and misses it with probability $1-1/n$. The draws are independent, so
\[
\Pr(i\ \text{not in sample})=\Bigl(1-\tfrac{1}{n}\Bigr)^{n}
\quad\Longrightarrow\quad
\Pr(i\ \text{in sample})=1-\Bigl(1-\tfrac{1}{n}\Bigr)^{n}.
\]
For $n=5$: $1-(4/5)^5=1-0.32768=0.672$.
As $n\to\infty$: $(1-1/n)^n\to e^{-1}$, so the probability tends to
$1-e^{-1}\approx 1-0.368=0.632$ --- a bootstrap sample contains about $63\,\%$ of the
distinct original observations, even for large $n$.

\medskip
\textbf{(d) Bootstrap standard error.}
$\bar{\alpha}^{*}=(0.4+0.5+0.6)/3=0.5$.
\[
\widehat{\operatorname{SE}}_B(\hat{\alpha})
=\sqrt{\frac{(0.4-0.5)^2+(0.5-0.5)^2+(0.6-0.5)^2}{3-1}}
=\sqrt{\frac{0.01+0+0.01}{2}}
=\sqrt{0.01}=0.100 .
\]
This estimates the standard deviation of $\hat{\alpha}$ over repeated samples from the
population, i.e.\ how much the estimated weight would typically vary if the data could be
collected anew. The bootstrap is useful because it replaces analytic derivations by
resampling from the observed data: it needs no formula for
$\operatorname{SE}(\hat{\alpha})$ and no new data, and it works for almost any statistic
--- including complicated ones for which no closed-form standard error exists.
\end{solutionbox}
\fi

% ============================================================
\problem{5: Model selection with $C_p$ and BIC}{12}

Best subset selection on a data set with $n=100$ observations produced, for each model
size $d$ (number of predictors), the best model and its residual sum of squares:

\begin{center}
\begin{tabular}{lcccc}
\toprule
Number of predictors $d$ & 1 & 2 & 3 & 4\\
\midrule
$\operatorname{RSS}_d$ & 500 & 444 & 428 & 424\\
\bottomrule
\end{tabular}
\end{center}

The noise variance is estimated from the full model as $\hat{\sigma}^2=4$.
Use the criteria
\[
C_p=\frac{1}{n}\bigl(\operatorname{RSS}+2\,d\,\hat{\sigma}^2\bigr),
\qquad
\operatorname{BIC}=\frac{1}{n}\bigl(\operatorname{RSS}+\log(n)\,d\,\hat{\sigma}^2\bigr),
\qquad \log(100)\approx 4.605 .
\]

\begin{enumerate}[label=(\alph*),leftmargin=7mm,itemsep=2.5mm]
  \item \pp{4} Compute $C_p$ for all four models (3 decimals).
  \item \pp{4} Compute BIC for all four models (3 decimals).
  \item \pp{2} Which model size does each criterion select? Explain the reason for the
    disagreement by comparing the two penalty terms.
  \item \pp{2} Why can the training RSS (or equivalently the training $R^2$) alone never be
    used to select $d$? What does RSS necessarily do as $d$ grows?
\end{enumerate}

\ifdefined\withsolutions
\begin{solutionbox}
\textbf{(a) $C_p$.} Penalty per predictor: $2\hat{\sigma}^2=8$; $C_p=(\operatorname{RSS}_d+8d)/100$:
\begin{center}
\begin{tabular}{lcccc}
\toprule
$d$ & 1 & 2 & 3 & 4\\
\midrule
$\operatorname{RSS}_d+8d$ & 508 & 460 & 452 & 456\\
$C_p$ & 5.080 & 4.600 & \textbf{4.520} & 4.560\\
\bottomrule
\end{tabular}
\end{center}

\medskip
\textbf{(b) BIC.} Penalty per predictor: $\log(100)\,\hat{\sigma}^2=4.605\cdot 4=18.421$;
$\operatorname{BIC}=(\operatorname{RSS}_d+18.421\,d)/100$:
\begin{center}
\begin{tabular}{lcccc}
\toprule
$d$ & 1 & 2 & 3 & 4\\
\midrule
$\operatorname{RSS}_d+18.421\,d$ & 518.421 & 480.841 & 483.262 & 497.683\\
BIC & 5.184 & \textbf{4.808} & 4.833 & 4.977\\
\bottomrule
\end{tabular}
\end{center}

\medskip
\textbf{(c) Selected models.}
$C_p$ is minimised at $d=3$ (4.520); BIC is minimised at $d=2$ (4.808).
Reason: both criteria are RSS plus a penalty for model size, but BIC charges
$\log(n)\hat{\sigma}^2=18.421$ per predictor whereas $C_p$ charges only
$2\hat{\sigma}^2=8$; since $\log(100)=4.605>2$, BIC penalises additional predictors much
more heavily and therefore tends to select \emph{smaller} models. Here the RSS drop from
$d=2$ to $d=3$ (16) exceeds the $C_p$ penalty (8) but not the BIC penalty (18.421).

\medskip
\textbf{(d) Why RSS alone cannot choose $d$.}
For nested model searches, adding a predictor can never increase the training RSS: RSS
decreases \emph{monotonically} in $d$ (and training $R^2$ increases monotonically).
RSS-based selection would therefore always pick the largest model, $d=4$, regardless of
whether the extra predictors carry real signal --- training error is an ever more
optimistic estimate of test error as flexibility grows. Hence one must either penalise
model size ($C_p$, AIC, BIC, adjusted $R^2$) or estimate the test error directly
(validation set, cross-validation).
\end{solutionbox}
\fi

% ============================================================
\problem{6: Ridge regression and the lasso}{15}

Consider a linear model with response $y$ and $p$ standardised predictors
$X_1,\dots,X_p$. Part (e) also draws on the bias--variance ideas of Chapter~2 and the
least-squares fit of Chapter~3.

\begin{enumerate}[label=(\alph*),leftmargin=7mm,itemsep=2.5mm]
  \item \pp{4} Write down the optimisation objectives of \emph{ridge regression} and of
    the \emph{lasso} (RSS plus the respective penalty, with tuning parameter
    $\lambda\ge 0$). Which coefficient is conventionally \emph{not} penalised?
  \item \pp{3} Describe the behaviour of the two estimators in the limits
    $\lambda\to 0$ and $\lambda\to\infty$. Which method from Chapter~3 is recovered as
    $\lambda\to 0$?
  \item \pp{3} Explain --- in words, using the constraint-region geometry for $p=2$ --- why
    the lasso can set coefficients \emph{exactly} to zero while ridge regression
    typically does not.
  \item \pp{2} Why should the predictors be \emph{standardised} before ridge or lasso is
    applied, when this is irrelevant for ordinary least squares?
  \item \pp{3} A ridge regression was fitted for two values of the tuning parameter; the
    table shows the standardised coefficient estimates. A 10-fold cross-validation over a
    grid of $\lambda$ values gave the CV errors on the right.
    \begin{center}
    \begin{tabular}{lrr}
    \toprule
     & $\lambda=0.1$ & $\lambda=100$\\
    \midrule
    $\hat{\beta}_1$ & $0.92$ & $0.31$\\
    $\hat{\beta}_2$ & $0.85$ & $0.29$\\
    $\hat{\beta}_3$ & $-0.41$ & $-0.12$\\
    $\hat{\beta}_4$ & $0.20$ & $0.06$\\
    \bottomrule
    \end{tabular}
    \hspace{12mm}
    \begin{tabular}{lcccc}
    \toprule
    $\lambda$ & 0.1 & 1 & 10 & 100\\
    \midrule
    CV error & 12.8 & 11.9 & 10.6 & 13.4\\
    \bottomrule
    \end{tabular}
    \end{center}
    Which of the two fits shown in the coefficient table has the lower bias and which the
    lower variance, and why? Which $\lambda$ from the CV table should be chosen for the
    final model?
\end{enumerate}

\ifdefined\withsolutions
\begin{solutionbox}
\textbf{(a) Objectives.} Ridge regression minimises
\[
\sum_{i=1}^{n}\Bigl(y_i-\beta_0-\sum_{j=1}^{p}\beta_j x_{ij}\Bigr)^{2}
+\lambda\sum_{j=1}^{p}\beta_j^{2}
=\operatorname{RSS}+\lambda\|\beta\|_2^2 ,
\]
the lasso minimises
\[
\sum_{i=1}^{n}\Bigl(y_i-\beta_0-\sum_{j=1}^{p}\beta_j x_{ij}\Bigr)^{2}
+\lambda\sum_{j=1}^{p}|\beta_j|
=\operatorname{RSS}+\lambda\|\beta\|_1 .
\]
The intercept $\beta_0$ is \emph{not} penalised (shrinking the overall level of $y$ would
make no sense).

\medskip
\textbf{(b) Limiting behaviour.}
$\lambda\to 0$: the penalty vanishes and both estimators approach the \emph{ordinary
least-squares} solution of Chapter~3.
$\lambda\to\infty$: the penalty dominates and all penalised coefficients are forced to
zero --- the null model containing only the intercept. Ridge shrinks the coefficients
\emph{towards} zero (exactly zero only in the limit), whereas the lasso sets more and more
coefficients \emph{exactly} to zero as $\lambda$ grows. In between, $\lambda$ traces out a
whole path of fits whose flexibility decreases as $\lambda$ increases.

\medskip
\textbf{(c) Geometry of the penalties ($p=2$).}
The penalised problems are equivalent to minimising the RSS subject to a budget
constraint: $|\beta_1|+|\beta_2|\le s$ (lasso) or $\beta_1^2+\beta_2^2\le s$ (ridge).
The lasso region is a \emph{diamond} with corners on the coordinate axes; the ridge region
is a \emph{disk} with no corners. The RSS is minimised over the constraint region where
its elliptical contours around the OLS solution first touch the region. Ellipses often hit
the diamond exactly at a \emph{corner} --- and at a corner one coefficient is exactly zero
(sparsity, variable selection). The circle has no corners, so the touching point generically
has both coordinates non-zero: ridge shrinks but keeps all $p$ predictors.

\medskip
\textbf{(d) Standardisation.}
OLS is scale-equivariant: rescaling $X_j$ (e.g.\ EUR $\to$ thousand EUR) simply rescales
$\hat{\beta}_j$, leaving fits unchanged. The penalties, however, treat all coefficients
symmetrically: $\lambda\sum\beta_j^2$ resp.\ $\lambda\sum|\beta_j|$ depend on the
\emph{magnitudes} of the coefficients and hence on the units of the predictors --- a
predictor measured in small units has a large coefficient and would be penalised
(shrunk) more strongly for purely cosmetic reasons. Standardising each predictor to
standard deviation one makes the penalty act comparably on all predictors and the solution
independent of the original units.

\medskip
\textbf{(e) Bias--variance and choice of $\lambda$.}
$\lambda=0.1$: almost no shrinkage --- the fit is close to OLS, i.e.\ more flexible:
\emph{lower bias, higher variance}. $\lambda=100$: strong shrinkage (coefficients roughly
one third of their former size, e.g.\ $0.92\to 0.31$): a more rigid fit with
\emph{higher bias but lower variance} --- the bias--variance trade-off of Chapter~2
governed by a single knob $\lambda$.
From the CV table, the estimated test error is minimised at $\lambda=10$
(CV error $10.6<11.9<12.8<13.4$), so \emph{$\lambda=10$} should be chosen --- an
intermediate amount of shrinkage; neither of the two fits shown in the coefficient table
is optimal.
\end{solutionbox}
\fi

\vspace{6mm}
\begin{center}
\rule{0.35\textwidth}{0.4pt}\\[1mm]
\textit{End of the examination --- good luck!}
\end{center}

\end{document}
"
% ============================================================
\documentclass[11pt,a4paper]{article}
\usepackage[utf8]{inputenc}
\usepackage[T1]{fontenc}
\usepackage[english]{babel}
\usepackage[margin=2.2cm]{geometry}
\usepackage{amsmath,amssymb}
\usepackage{booktabs,array}
\usepackage{enumitem}
\usepackage{xcolor}
\usepackage{tcolorbox}
\tcbuselibrary{skins,breakable}

\setlength{\parindent}{0pt}

% Teal solution box (only shown when \withsolutions is defined)
\newtcolorbox{solutionbox}{enhanced,breakable,colback=teal!5,colframe=teal!55!black,
  boxrule=0pt,leftrule=2.5pt,arc=2pt,fonttitle=\bfseries,title=Solution,
  left=2mm,right=2mm,top=1mm,bottom=1.5mm}

\newcommand{\problem}[2]{\par\vspace{7mm}\noindent{\large\textbf{Problem #1}}%
  \hfill{\large\textbf{(#2 points)}}\par\nopagebreak\vspace{0.6mm}%
  \noindent\rule{\textwidth}{0.6pt}\par\nopagebreak\vspace{2.5mm}}
\newcommand{\pp}[1]{\textbf{[#1~P]}\enspace}

\begin{document}

% ------------------------------------------------------------
% Header block
% ------------------------------------------------------------
\begin{center}
  {\Large\textbf{HSBI --- Bielefeld University of Applied Sciences and Arts}}\\[1.2mm]
  {\large Quantitative Research Methods}\\[1mm]
  {\normalsize Summer Semester 2026 \quad$\cdot$\quad Examiner: Prof.\ Dr.\ Christoph Weisser}\\[4.5mm]
  {\LARGE\textbf{Mock Examination 2 (Lectures 5--8)}}\\[1.2mm]
  {\normalsize ISLP Chapters 4--6: Classification $\cdot$ Resampling Methods $\cdot$ Model Selection and Regularisation}%
  \ifdefined\withsolutions\\[2.2mm]{\large\textcolor{teal!55!black}{\textbf{--- Version with detailed solutions ---}}}\fi
\end{center}

\vspace{2.5mm}
\begin{center}
\begin{tabular}{@{}ll@{}}
  \textbf{Duration:} & \textbf{90 minutes}\\
  \textbf{Total credit:} & \textbf{90 points}\\
  \textbf{Allowed aids:} & non-programmable calculator; one handwritten A4 sheet of notes (both sides)\\
\end{tabular}
\end{center}

\vspace{3mm}
Name: \rule[-1mm]{0.38\textwidth}{0.4pt}\hfill Matriculation no.: \rule[-1mm]{0.24\textwidth}{0.4pt}

\vspace{4.5mm}
\begin{center}
\renewcommand{\arraystretch}{1.35}
\begin{tabular}{|l|c|c|c|c|c|c|c|}
\hline
\textbf{Problem} & \textbf{P1} & \textbf{P2} & \textbf{P3} & \textbf{P4} & \textbf{P5} & \textbf{P6} & \textbf{Total}\\\hline
Maximum points & 18 & 12 & 15 & 18 & 12 & 15 & \textbf{90}\\\hline
Achieved points & \hspace{9mm} & \hspace{9mm} & \hspace{9mm} & \hspace{9mm} & \hspace{9mm} & \hspace{9mm} & \hspace{11mm}\\\hline
\end{tabular}
\end{center}

\vspace{2.5mm}
\textbf{Instructions}
\begin{itemize}[nosep,leftmargin=5.5mm]
  \item Justify all answers. Numerical results without a derivation or a brief reasoning receive no credit.
  \item Round final numerical results to 3 decimal places unless stated otherwise. Small deviations caused by carrying rounded intermediate results are accepted.
  \item The number of points per problem roughly equals the recommended working time in minutes.
\end{itemize}
\vspace{1mm}\noindent\rule{\textwidth}{0.6pt}

% ============================================================
\problem{1: Logistic regression by hand}{18}

A bank models the probability that a credit-card customer \emph{defaults} ($Y=1$) as a
function of the monthly card balance $x$ (in EUR; balances in this portfolio range from
0 to 400 EUR). A logistic regression fitted by maximum likelihood gave
\[
\hat{p}(x)=\Pr(Y=1\mid X=x)
  =\frac{e^{\hat{\beta}_0+\hat{\beta}_1 x}}{1+e^{\hat{\beta}_0+\hat{\beta}_1 x}},
\qquad
\hat{\beta}_0=-4,\quad \hat{\beta}_1=0.02 .
\]

\begin{enumerate}[label=(\alph*),leftmargin=7mm,itemsep=2.5mm]
  \item \pp{4} Compute the predicted default probability $\hat{p}(x)$ for a customer with
    balance $x=100$ EUR and for a customer with balance $x=250$ EUR.
  \item \pp{4} For the same two customers, compute the \emph{odds} of default
    $\hat{p}(x)/(1-\hat{p}(x))$ and the \emph{log-odds} (logit). Interpret the odds of the
    customer with $x=250$ in one sentence.
  \item \pp{3} Determine the balance $x^{*}$ at which the predicted default probability
    equals exactly $0.5$. Show your derivation.
  \item \pp{4} Interpret the estimated slope on the odds scale: by which \emph{factor} do the
    estimated odds of default change (i) per additional EUR of balance and (ii) per
    additional 100 EUR of balance? Why is the statement ``$\hat{\beta}_1=0.02$ means the
    default \emph{probability} rises by 0.02 per EUR'' wrong?
  \item \pp{3} Give two reasons why an ordinary linear regression of the 0/1 response $Y$
    on $x$ would be inappropriate here, even though it can be computed mechanically.
\end{enumerate}

\ifdefined\withsolutions
\begin{solutionbox}
\textbf{(a) Predicted probabilities.} The linear predictor (log-odds) is
$\hat{\eta}(x)=\hat{\beta}_0+\hat{\beta}_1x=-4+0.02x$.
\[
x=100:\ \hat{\eta}=-4+2=-2,\qquad
\hat{p}(100)=\frac{e^{-2}}{1+e^{-2}}=\frac{0.1353}{1.1353}=0.119 .
\]
\[
x=250:\ \hat{\eta}=-4+5=1,\qquad
\hat{p}(250)=\frac{e^{1}}{1+e^{1}}=\frac{2.7183}{3.7183}=0.731 .
\]

\medskip
\textbf{(b) Odds and log-odds.} Since
$\hat{p}/(1-\hat{p})=e^{\hat{\eta}}$ and $\log\big(\hat{p}/(1-\hat{p})\big)=\hat{\eta}$:
\begin{center}
\begin{tabular}{lcc}
\toprule
 & $x=100$ & $x=250$\\
\midrule
log-odds $\hat{\eta}(x)$ & $-2.000$ & $1.000$\\
odds $e^{\hat{\eta}(x)}$ & $e^{-2}=0.135$ & $e^{1}=2.718$\\
\bottomrule
\end{tabular}
\end{center}
Check via (a): $0.119/0.881=0.135$ and $0.731/0.269=2.718$. \checkmark\
\emph{Interpretation for $x=250$:} the odds are about $2.72$ to 1, i.e.\ default is
estimated to be about $2.7$ times as probable as non-default for this customer.

\medskip
\textbf{(c) Balance with $\hat{p}=0.5$.}
$\hat{p}(x^{*})=0.5$ $\iff$ odds $=1$ $\iff$ log-odds $=0$:
\[
-4+0.02\,x^{*}=0
\quad\Longrightarrow\quad
x^{*}=\frac{4}{0.02}=200\ \text{EUR}.
\]
Below 200 EUR the model predicts ``no default'' as more likely, above 200 EUR ``default''.

\medskip
\textbf{(d) Odds-scale interpretation.}
Per additional EUR the odds are multiplied by
$e^{\hat{\beta}_1}=e^{0.02}=1.020$, i.e.\ they increase by about $2.0\,\%$.
Per additional 100 EUR they are multiplied by
$e^{100\hat{\beta}_1}=e^{2}=7.389$ (an increase of about $639\,\%$).
The quoted statement is wrong because $\hat{\beta}_1$ is the change in the
\emph{log-odds} per EUR, not in the probability: the probability change
$\partial\hat{p}/\partial x=\hat{\beta}_1\hat{p}(x)\bigl(1-\hat{p}(x)\bigr)$ depends on $x$
--- the S-shaped curve is steepest near $\hat{p}=0.5$ and nearly flat in the tails, so there
is no constant ``probability effect per EUR''.

\medskip
\textbf{(e) Why not linear regression?}
(1) A straight line $\hat{\beta}_0+\hat{\beta}_1x$ is unbounded, so for small or large
balances it produces ``probabilities'' below 0 or above 1, which are not sensible
estimates of $\Pr(Y=1\mid x)$; the logistic function is guaranteed to stay in $[0,1]$.
(2) The 0/1 response violates the linear-model error assumptions: the errors are not
normal and their variance $p(x)(1-p(x))$ depends on $x$ (heteroscedasticity), so the usual
inference is unreliable. (Also, with more than two classes a numeric coding of $Y$ would
impose an artificial ordering and equal spacing between classes.)
\end{solutionbox}
\fi

% ============================================================
\problem{2: Discriminant analysis and naive Bayes --- concepts}{12}

\begin{enumerate}[label=(\alph*),leftmargin=7mm,itemsep=2.5mm]
  \item \pp{4} Linear discriminant analysis (LDA) and quadratic discriminant analysis (QDA)
    both model the class-conditional densities $f_k(x)$ of the predictors as multivariate
    normal. State precisely the assumption about the covariance matrices that distinguishes
    the two methods, and give the number of covariance \emph{matrices} each method estimates
    when there are $K$ classes.
  \item \pp{3} Using the bias--variance trade-off, explain in which situations LDA tends to
    outperform QDA on test data, and in which situations QDA tends to win. Mention the role
    of the sample size $n$ and of the true class covariances.
  \item \pp{3} State the key assumption of the \emph{naive Bayes} classifier about the
    predictors within each class, and explain why this assumption can improve test
    performance when the number of predictors $p$ is large relative to $n$ --- even though
    the assumption is rarely exactly true.
  \item \pp{2} Which of the three classifiers (LDA, QDA, naive Bayes) produce \emph{linear}
    decision boundaries in $x$, and which produce \emph{quadratic} or more general
    non-linear boundaries?
\end{enumerate}

\ifdefined\withsolutions
\begin{solutionbox}
\textbf{(a) Covariance assumptions.}
Both methods assume $X\mid Y=k\sim\mathcal{N}(\mu_k,\Sigma_k)$.
\emph{LDA} additionally assumes a covariance matrix that is \emph{shared} by all classes:
$\Sigma_1=\dots=\Sigma_K=\Sigma$; it therefore estimates \emph{one} common $p\times p$
covariance matrix. \emph{QDA} allows each class its \emph{own} covariance matrix
$\Sigma_k$ and estimates $K$ covariance matrices (one per class), i.e.\ roughly
$K\,p(p+1)/2$ covariance parameters instead of $p(p+1)/2$.

\medskip
\textbf{(b) When does which method win?}
QDA is the more flexible method (many more parameters): \emph{low bias, high variance};
LDA is more restrictive: \emph{higher bias if the $\Sigma_k$ truly differ, but much lower
variance}. Hence LDA tends to win when $n$ is small (or $p$ large), so that reducing
variance is crucial, and when the true class covariances are (approximately) equal ---
then LDA's assumption costs little bias. QDA tends to win when $n$ is large, so its extra
parameters can be estimated reliably, and when the class covariances clearly differ, so
LDA's shared-$\Sigma$ assumption creates substantial bias.

\medskip
\textbf{(c) Naive Bayes.}
Assumption: \emph{within each class the $p$ predictors are independent}, i.e.\ the joint
density factorises as $f_k(x)=f_{k1}(x_1)\cdot f_{k2}(x_2)\cdots f_{kp}(x_p)$.
Instead of a full $p$-dimensional joint density (which would require estimating
associations between all predictors, e.g.\ $p(p+1)/2$ covariance parameters per class),
only $p$ one-dimensional marginal densities per class must be estimated. This drastic
reduction in the number of parameters lowers the \emph{variance} of the classifier.
When $p$ is large relative to $n$, the variance saving usually outweighs the bias
introduced by (falsely) ignoring within-class dependence --- another instance of the
bias--variance trade-off.

\medskip
\textbf{(d) Shape of the decision boundaries.}
\emph{LDA:} linear boundaries (the discriminant functions are linear in $x$; the quadratic
term cancels because $\Sigma$ is shared). \emph{QDA:} quadratic boundaries (each class has
its own quadratic form). \emph{Naive Bayes:} generally non-linear boundaries of flexible
shape --- the log-odds are an \emph{additive} function $\sum_j g_j(x_j)$ of the individual
predictors (no interaction terms); only in special cases (e.g.\ Gaussian marginals with
class-shared variances) does it reduce to a linear rule.
\end{solutionbox}
\fi

% ============================================================
\problem{3: Confusion matrix and ROC curve}{15}

A logistic regression classifier predicts credit-card default on an evaluation set of
$n=1{,}000$ customers, using the default threshold: predict ``default'' if
$\hat{p}(\text{default}\mid x)>0.5$. The result is the following confusion matrix:

\begin{center}
\renewcommand{\arraystretch}{1.15}
\begin{tabular}{llcc|c}
 & & \multicolumn{2}{c|}{\textbf{True status}} & \\
 & & Default & No default & Total\\
\cmidrule{2-5}
\textbf{Predicted} & Default    & 60 & 45  & 105\\
                   & No default & 40 & 855 & 895\\
\cmidrule{2-5}
 & Total & 100 & 900 & 1{,}000\\
\end{tabular}
\end{center}

\begin{enumerate}[label=(\alph*),leftmargin=7mm,itemsep=2.5mm]
  \item \pp{7} Identify TP, FP, FN and TN, and compute (i) the overall accuracy and the
    overall error rate, (ii) the sensitivity (recall, true positive rate), (iii) the
    specificity, (iv) the precision, and (v) the false positive rate. Interpret sensitivity
    and precision in one sentence each in the context of this application.
  \item \pp{4} The bank now lowers the classification threshold from $0.5$ to $0.2$
    (predict ``default'' if $\hat{p}>0.2$). Explain --- without further computation --- how
    sensitivity and specificity will change, and give a business reason why the bank might
    deliberately accept this trade-off.
  \item \pp{4} The ROC curve plots which two quantities against each other, as what is
    varied? What does an area under the curve (AUC) of $0.5$ mean, and what does an AUC of
    $1$ mean?
\end{enumerate}

\ifdefined\withsolutions
\begin{solutionbox}
\textbf{(a) Metrics.} With ``default'' as the positive class:
$\text{TP}=60$, $\text{FP}=45$, $\text{FN}=40$, $\text{TN}=855$.
\begin{align*}
\text{(i) accuracy} &= \frac{\text{TP}+\text{TN}}{n}=\frac{60+855}{1000}=\frac{915}{1000}=0.915,
\qquad \text{error rate}=1-0.915=0.085,\\[0.5mm]
\text{(ii) sensitivity} &= \frac{\text{TP}}{\text{TP}+\text{FN}}=\frac{60}{100}=0.600,\\[0.5mm]
\text{(iii) specificity} &= \frac{\text{TN}}{\text{TN}+\text{FP}}=\frac{855}{900}=0.950,\\[0.5mm]
\text{(iv) precision} &= \frac{\text{TP}}{\text{TP}+\text{FP}}=\frac{60}{105}=0.571,\\[0.5mm]
\text{(v) FPR} &= \frac{\text{FP}}{\text{FP}+\text{TN}}=\frac{45}{900}=0.050=1-\text{specificity}.
\end{align*}
\emph{Sensitivity:} the classifier detects $60.0\,\%$ of the customers who actually
default --- $40\,\%$ of true defaulters are missed.
\emph{Precision:} of the customers flagged as defaulters, only $57.1\,\%$ actually
default; the rest are false alarms.

\medskip
\textbf{(b) Lowering the threshold to 0.2.}
More customers exceed the lower cut-off, so more are predicted ``default'': every true
defaulter with $0.2<\hat{p}\le 0.5$ is now caught, so \emph{sensitivity increases}
(FN $\downarrow$); at the same time more non-defaulters are flagged, so
\emph{specificity decreases} (FP $\uparrow$, FPR $\uparrow$). (The overall error rate will
typically increase, since non-defaulters are the large majority.)
\emph{Business reason:} the costs are asymmetric --- a missed defaulter (FN, a credit
loss) is far more expensive than a false alarm (FP, e.g.\ an unnecessary credit-line
review or customer call). Trading specificity for sensitivity can therefore minimise the
\emph{expected cost} even if the raw error rate rises.

\medskip
\textbf{(c) ROC and AUC.}
The ROC curve plots the \emph{true positive rate} (sensitivity) on the vertical axis
against the \emph{false positive rate} ($1-$specificity) on the horizontal axis, as the
classification \emph{threshold} is varied over all values from 1 down to 0.
$\text{AUC}=0.5$: performance equal to the diagonal, i.e.\ no better than random guessing
--- the score contains no usable ranking information. $\text{AUC}=1$: perfect
classifier --- some threshold separates the two classes without error
(sensitivity $=1$ and specificity $=1$ simultaneously). Equivalently, the AUC is the
probability that a randomly chosen defaulter receives a higher score $\hat{p}$ than a
randomly chosen non-defaulter.
\end{solutionbox}
\fi

% ============================================================
\problem{4: Resampling methods}{18}

\begin{enumerate}[label=(\alph*),leftmargin=7mm,itemsep=2.5mm]
  \item \pp{4} A data set has $n=200$ observations. To estimate the test error of a fixed
    model specification, state how many times the model must be \emph{fitted} for (i) the
    validation-set approach (one 50/50 split), (ii) 10-fold cross-validation, and (iii)
    leave-one-out cross-validation (LOOCV), with a one-sentence justification each.
  \item \pp{4} Compare LOOCV and 10-fold CV as estimators of the true test error with
    respect to \emph{bias} and \emph{variance}. Which method is preferred in practice, and
    why?
  \item \pp{5} A bootstrap sample of size $n$ is drawn with replacement from $n$
    observations. Show that the probability that a fixed observation $i$ is
    \emph{contained} in the bootstrap sample equals $1-(1-1/n)^n$. Evaluate this
    probability for $n=5$, and state the limiting value as $n\to\infty$
    (use $e^{-1}\approx 0.368$).
  \item \pp{5} To quantify the uncertainty of an estimated portfolio weight
    $\hat{\alpha}$, an analyst draws $B=3$ bootstrap samples (in practice $B\approx 1{,}000$;
    $B=3$ is used here only to keep the arithmetic manageable) and obtains
    \[
    \hat{\alpha}^{*1}=0.4,\qquad \hat{\alpha}^{*2}=0.5,\qquad \hat{\alpha}^{*3}=0.6 .
    \]
    Compute the bootstrap estimate of the standard error,
    \[
    \widehat{\operatorname{SE}}_B(\hat{\alpha})
      =\sqrt{\frac{1}{B-1}\sum_{r=1}^{B}\Bigl(\hat{\alpha}^{*r}-\bar{\alpha}^{*}\Bigr)^{2}},
    \qquad
    \bar{\alpha}^{*}=\frac{1}{B}\sum_{r=1}^{B}\hat{\alpha}^{*r},
    \]
    and explain in one sentence what this number estimates. Why is the bootstrap useful
    here although a formula for $\operatorname{SE}(\hat{\alpha})$ may be hard to derive?
\end{enumerate}

\ifdefined\withsolutions
\begin{solutionbox}
\textbf{(a) Number of model fits at $n=200$.}
(i) \emph{Validation set:} \textbf{1} fit --- the model is trained once on the training
half (100 observations) and evaluated on the held-out half.
(ii) \emph{10-fold CV:} \textbf{10} fits --- each of the 10 folds (20 observations) is
held out once, and the model is refitted on the remaining 180 observations each time.
(iii) \emph{LOOCV:} \textbf{200} fits --- each single observation is held out once and
the model is refitted on the other 199. (LOOCV $=$ $n$-fold CV; for least squares a
shortcut formula avoids the refits, but generically $n$ fits are needed.)

\medskip
\textbf{(b) Bias and variance of the CV estimators.}
\emph{Bias:} LOOCV trains on $n-1$ observations, almost the full sample, so it is nearly
unbiased for the true test error; 10-fold CV trains on $0.9\,n$ and therefore slightly
\emph{overestimates} the test error (a bit more bias). \emph{Variance:} LOOCV has the
\emph{higher} variance: the $n$ fitted models are trained on almost identical data sets,
so their errors are strongly positively correlated, and the mean of highly correlated
quantities has high variance; the 10 folds of 10-fold CV overlap less, giving less
correlated fits and a lower-variance estimate. \emph{In practice} $k=5$ or $k=10$ is
preferred: a good bias--variance compromise and much cheaper computationally.

\medskip
\textbf{(c) Bootstrap inclusion probability.}
Each of the $n$ draws is made with replacement, so each draw selects observation $i$ with
probability $1/n$ and misses it with probability $1-1/n$. The draws are independent, so
\[
\Pr(i\ \text{not in sample})=\Bigl(1-\tfrac{1}{n}\Bigr)^{n}
\quad\Longrightarrow\quad
\Pr(i\ \text{in sample})=1-\Bigl(1-\tfrac{1}{n}\Bigr)^{n}.
\]
For $n=5$: $1-(4/5)^5=1-0.32768=0.672$.
As $n\to\infty$: $(1-1/n)^n\to e^{-1}$, so the probability tends to
$1-e^{-1}\approx 1-0.368=0.632$ --- a bootstrap sample contains about $63\,\%$ of the
distinct original observations, even for large $n$.

\medskip
\textbf{(d) Bootstrap standard error.}
$\bar{\alpha}^{*}=(0.4+0.5+0.6)/3=0.5$.
\[
\widehat{\operatorname{SE}}_B(\hat{\alpha})
=\sqrt{\frac{(0.4-0.5)^2+(0.5-0.5)^2+(0.6-0.5)^2}{3-1}}
=\sqrt{\frac{0.01+0+0.01}{2}}
=\sqrt{0.01}=0.100 .
\]
This estimates the standard deviation of $\hat{\alpha}$ over repeated samples from the
population, i.e.\ how much the estimated weight would typically vary if the data could be
collected anew. The bootstrap is useful because it replaces analytic derivations by
resampling from the observed data: it needs no formula for
$\operatorname{SE}(\hat{\alpha})$ and no new data, and it works for almost any statistic
--- including complicated ones for which no closed-form standard error exists.
\end{solutionbox}
\fi

% ============================================================
\problem{5: Model selection with $C_p$ and BIC}{12}

Best subset selection on a data set with $n=100$ observations produced, for each model
size $d$ (number of predictors), the best model and its residual sum of squares:

\begin{center}
\begin{tabular}{lcccc}
\toprule
Number of predictors $d$ & 1 & 2 & 3 & 4\\
\midrule
$\operatorname{RSS}_d$ & 500 & 444 & 428 & 424\\
\bottomrule
\end{tabular}
\end{center}

The noise variance is estimated from the full model as $\hat{\sigma}^2=4$.
Use the criteria
\[
C_p=\frac{1}{n}\bigl(\operatorname{RSS}+2\,d\,\hat{\sigma}^2\bigr),
\qquad
\operatorname{BIC}=\frac{1}{n}\bigl(\operatorname{RSS}+\log(n)\,d\,\hat{\sigma}^2\bigr),
\qquad \log(100)\approx 4.605 .
\]

\begin{enumerate}[label=(\alph*),leftmargin=7mm,itemsep=2.5mm]
  \item \pp{4} Compute $C_p$ for all four models (3 decimals).
  \item \pp{4} Compute BIC for all four models (3 decimals).
  \item \pp{2} Which model size does each criterion select? Explain the reason for the
    disagreement by comparing the two penalty terms.
  \item \pp{2} Why can the training RSS (or equivalently the training $R^2$) alone never be
    used to select $d$? What does RSS necessarily do as $d$ grows?
\end{enumerate}

\ifdefined\withsolutions
\begin{solutionbox}
\textbf{(a) $C_p$.} Penalty per predictor: $2\hat{\sigma}^2=8$; $C_p=(\operatorname{RSS}_d+8d)/100$:
\begin{center}
\begin{tabular}{lcccc}
\toprule
$d$ & 1 & 2 & 3 & 4\\
\midrule
$\operatorname{RSS}_d+8d$ & 508 & 460 & 452 & 456\\
$C_p$ & 5.080 & 4.600 & \textbf{4.520} & 4.560\\
\bottomrule
\end{tabular}
\end{center}

\medskip
\textbf{(b) BIC.} Penalty per predictor: $\log(100)\,\hat{\sigma}^2=4.605\cdot 4=18.421$;
$\operatorname{BIC}=(\operatorname{RSS}_d+18.421\,d)/100$:
\begin{center}
\begin{tabular}{lcccc}
\toprule
$d$ & 1 & 2 & 3 & 4\\
\midrule
$\operatorname{RSS}_d+18.421\,d$ & 518.421 & 480.841 & 483.262 & 497.683\\
BIC & 5.184 & \textbf{4.808} & 4.833 & 4.977\\
\bottomrule
\end{tabular}
\end{center}

\medskip
\textbf{(c) Selected models.}
$C_p$ is minimised at $d=3$ (4.520); BIC is minimised at $d=2$ (4.808).
Reason: both criteria are RSS plus a penalty for model size, but BIC charges
$\log(n)\hat{\sigma}^2=18.421$ per predictor whereas $C_p$ charges only
$2\hat{\sigma}^2=8$; since $\log(100)=4.605>2$, BIC penalises additional predictors much
more heavily and therefore tends to select \emph{smaller} models. Here the RSS drop from
$d=2$ to $d=3$ (16) exceeds the $C_p$ penalty (8) but not the BIC penalty (18.421).

\medskip
\textbf{(d) Why RSS alone cannot choose $d$.}
For nested model searches, adding a predictor can never increase the training RSS: RSS
decreases \emph{monotonically} in $d$ (and training $R^2$ increases monotonically).
RSS-based selection would therefore always pick the largest model, $d=4$, regardless of
whether the extra predictors carry real signal --- training error is an ever more
optimistic estimate of test error as flexibility grows. Hence one must either penalise
model size ($C_p$, AIC, BIC, adjusted $R^2$) or estimate the test error directly
(validation set, cross-validation).
\end{solutionbox}
\fi

% ============================================================
\problem{6: Ridge regression and the lasso}{15}

Consider a linear model with response $y$ and $p$ standardised predictors
$X_1,\dots,X_p$. Part (e) also draws on the bias--variance ideas of Chapter~2 and the
least-squares fit of Chapter~3.

\begin{enumerate}[label=(\alph*),leftmargin=7mm,itemsep=2.5mm]
  \item \pp{4} Write down the optimisation objectives of \emph{ridge regression} and of
    the \emph{lasso} (RSS plus the respective penalty, with tuning parameter
    $\lambda\ge 0$). Which coefficient is conventionally \emph{not} penalised?
  \item \pp{3} Describe the behaviour of the two estimators in the limits
    $\lambda\to 0$ and $\lambda\to\infty$. Which method from Chapter~3 is recovered as
    $\lambda\to 0$?
  \item \pp{3} Explain --- in words, using the constraint-region geometry for $p=2$ --- why
    the lasso can set coefficients \emph{exactly} to zero while ridge regression
    typically does not.
  \item \pp{2} Why should the predictors be \emph{standardised} before ridge or lasso is
    applied, when this is irrelevant for ordinary least squares?
  \item \pp{3} A ridge regression was fitted for two values of the tuning parameter; the
    table shows the standardised coefficient estimates. A 10-fold cross-validation over a
    grid of $\lambda$ values gave the CV errors on the right.
    \begin{center}
    \begin{tabular}{lrr}
    \toprule
     & $\lambda=0.1$ & $\lambda=100$\\
    \midrule
    $\hat{\beta}_1$ & $0.92$ & $0.31$\\
    $\hat{\beta}_2$ & $0.85$ & $0.29$\\
    $\hat{\beta}_3$ & $-0.41$ & $-0.12$\\
    $\hat{\beta}_4$ & $0.20$ & $0.06$\\
    \bottomrule
    \end{tabular}
    \hspace{12mm}
    \begin{tabular}{lcccc}
    \toprule
    $\lambda$ & 0.1 & 1 & 10 & 100\\
    \midrule
    CV error & 12.8 & 11.9 & 10.6 & 13.4\\
    \bottomrule
    \end{tabular}
    \end{center}
    Which of the two fits shown in the coefficient table has the lower bias and which the
    lower variance, and why? Which $\lambda$ from the CV table should be chosen for the
    final model?
\end{enumerate}

\ifdefined\withsolutions
\begin{solutionbox}
\textbf{(a) Objectives.} Ridge regression minimises
\[
\sum_{i=1}^{n}\Bigl(y_i-\beta_0-\sum_{j=1}^{p}\beta_j x_{ij}\Bigr)^{2}
+\lambda\sum_{j=1}^{p}\beta_j^{2}
=\operatorname{RSS}+\lambda\|\beta\|_2^2 ,
\]
the lasso minimises
\[
\sum_{i=1}^{n}\Bigl(y_i-\beta_0-\sum_{j=1}^{p}\beta_j x_{ij}\Bigr)^{2}
+\lambda\sum_{j=1}^{p}|\beta_j|
=\operatorname{RSS}+\lambda\|\beta\|_1 .
\]
The intercept $\beta_0$ is \emph{not} penalised (shrinking the overall level of $y$ would
make no sense).

\medskip
\textbf{(b) Limiting behaviour.}
$\lambda\to 0$: the penalty vanishes and both estimators approach the \emph{ordinary
least-squares} solution of Chapter~3.
$\lambda\to\infty$: the penalty dominates and all penalised coefficients are forced to
zero --- the null model containing only the intercept. Ridge shrinks the coefficients
\emph{towards} zero (exactly zero only in the limit), whereas the lasso sets more and more
coefficients \emph{exactly} to zero as $\lambda$ grows. In between, $\lambda$ traces out a
whole path of fits whose flexibility decreases as $\lambda$ increases.

\medskip
\textbf{(c) Geometry of the penalties ($p=2$).}
The penalised problems are equivalent to minimising the RSS subject to a budget
constraint: $|\beta_1|+|\beta_2|\le s$ (lasso) or $\beta_1^2+\beta_2^2\le s$ (ridge).
The lasso region is a \emph{diamond} with corners on the coordinate axes; the ridge region
is a \emph{disk} with no corners. The RSS is minimised over the constraint region where
its elliptical contours around the OLS solution first touch the region. Ellipses often hit
the diamond exactly at a \emph{corner} --- and at a corner one coefficient is exactly zero
(sparsity, variable selection). The circle has no corners, so the touching point generically
has both coordinates non-zero: ridge shrinks but keeps all $p$ predictors.

\medskip
\textbf{(d) Standardisation.}
OLS is scale-equivariant: rescaling $X_j$ (e.g.\ EUR $\to$ thousand EUR) simply rescales
$\hat{\beta}_j$, leaving fits unchanged. The penalties, however, treat all coefficients
symmetrically: $\lambda\sum\beta_j^2$ resp.\ $\lambda\sum|\beta_j|$ depend on the
\emph{magnitudes} of the coefficients and hence on the units of the predictors --- a
predictor measured in small units has a large coefficient and would be penalised
(shrunk) more strongly for purely cosmetic reasons. Standardising each predictor to
standard deviation one makes the penalty act comparably on all predictors and the solution
independent of the original units.

\medskip
\textbf{(e) Bias--variance and choice of $\lambda$.}
$\lambda=0.1$: almost no shrinkage --- the fit is close to OLS, i.e.\ more flexible:
\emph{lower bias, higher variance}. $\lambda=100$: strong shrinkage (coefficients roughly
one third of their former size, e.g.\ $0.92\to 0.31$): a more rigid fit with
\emph{higher bias but lower variance} --- the bias--variance trade-off of Chapter~2
governed by a single knob $\lambda$.
From the CV table, the estimated test error is minimised at $\lambda=10$
(CV error $10.6<11.9<12.8<13.4$), so \emph{$\lambda=10$} should be chosen --- an
intermediate amount of shrinkage; neither of the two fits shown in the coefficient table
is optimal.
\end{solutionbox}
\fi

\vspace{6mm}
\begin{center}
\rule{0.35\textwidth}{0.4pt}\\[1mm]
\textit{End of the examination --- good luck!}
\end{center}

\end{document}
"
% ============================================================
\documentclass[11pt,a4paper]{article}
\usepackage[utf8]{inputenc}
\usepackage[T1]{fontenc}
\usepackage[english]{babel}
\usepackage[margin=2.2cm]{geometry}
\usepackage{amsmath,amssymb}
\usepackage{booktabs,array}
\usepackage{enumitem}
\usepackage{xcolor}
\usepackage{tcolorbox}
\tcbuselibrary{skins,breakable}

\setlength{\parindent}{0pt}

% Teal solution box (only shown when \withsolutions is defined)
\newtcolorbox{solutionbox}{enhanced,breakable,colback=teal!5,colframe=teal!55!black,
  boxrule=0pt,leftrule=2.5pt,arc=2pt,fonttitle=\bfseries,title=Solution,
  left=2mm,right=2mm,top=1mm,bottom=1.5mm}

\newcommand{\problem}[2]{\par\vspace{7mm}\noindent{\large\textbf{Problem #1}}%
  \hfill{\large\textbf{(#2 points)}}\par\nopagebreak\vspace{0.6mm}%
  \noindent\rule{\textwidth}{0.6pt}\par\nopagebreak\vspace{2.5mm}}
\newcommand{\pp}[1]{\textbf{[#1~P]}\enspace}

\begin{document}

% ------------------------------------------------------------
% Header block
% ------------------------------------------------------------
\begin{center}
  {\Large\textbf{HSBI --- Bielefeld University of Applied Sciences and Arts}}\\[1.2mm]
  {\large Quantitative Research Methods}\\[1mm]
  {\normalsize Summer Semester 2026 \quad$\cdot$\quad Examiner: Prof.\ Dr.\ Christoph Weisser}\\[4.5mm]
  {\LARGE\textbf{Mock Examination 2 (Lectures 5--8)}}\\[1.2mm]
  {\normalsize ISLP Chapters 4--6: Classification $\cdot$ Resampling Methods $\cdot$ Model Selection and Regularisation}%
  \ifdefined\withsolutions\\[2.2mm]{\large\textcolor{teal!55!black}{\textbf{--- Version with detailed solutions ---}}}\fi
\end{center}

\vspace{2.5mm}
\begin{center}
\begin{tabular}{@{}ll@{}}
  \textbf{Duration:} & \textbf{90 minutes}\\
  \textbf{Total credit:} & \textbf{90 points}\\
  \textbf{Allowed aids:} & non-programmable calculator; one handwritten A4 sheet of notes (both sides)\\
\end{tabular}
\end{center}

\vspace{3mm}
Name: \rule[-1mm]{0.38\textwidth}{0.4pt}\hfill Matriculation no.: \rule[-1mm]{0.24\textwidth}{0.4pt}

\vspace{4.5mm}
\begin{center}
\renewcommand{\arraystretch}{1.35}
\begin{tabular}{|l|c|c|c|c|c|c|c|}
\hline
\textbf{Problem} & \textbf{P1} & \textbf{P2} & \textbf{P3} & \textbf{P4} & \textbf{P5} & \textbf{P6} & \textbf{Total}\\\hline
Maximum points & 18 & 12 & 15 & 18 & 12 & 15 & \textbf{90}\\\hline
Achieved points & \hspace{9mm} & \hspace{9mm} & \hspace{9mm} & \hspace{9mm} & \hspace{9mm} & \hspace{9mm} & \hspace{11mm}\\\hline
\end{tabular}
\end{center}

\vspace{2.5mm}
\textbf{Instructions}
\begin{itemize}[nosep,leftmargin=5.5mm]
  \item Justify all answers. Numerical results without a derivation or a brief reasoning receive no credit.
  \item Round final numerical results to 3 decimal places unless stated otherwise. Small deviations caused by carrying rounded intermediate results are accepted.
  \item The number of points per problem roughly equals the recommended working time in minutes.
\end{itemize}
\vspace{1mm}\noindent\rule{\textwidth}{0.6pt}

% ============================================================
\problem{1: Logistic regression by hand}{18}

A bank models the probability that a credit-card customer \emph{defaults} ($Y=1$) as a
function of the monthly card balance $x$ (in EUR; balances in this portfolio range from
0 to 400 EUR). A logistic regression fitted by maximum likelihood gave
\[
\hat{p}(x)=\Pr(Y=1\mid X=x)
  =\frac{e^{\hat{\beta}_0+\hat{\beta}_1 x}}{1+e^{\hat{\beta}_0+\hat{\beta}_1 x}},
\qquad
\hat{\beta}_0=-4,\quad \hat{\beta}_1=0.02 .
\]

\begin{enumerate}[label=(\alph*),leftmargin=7mm,itemsep=2.5mm]
  \item \pp{4} Compute the predicted default probability $\hat{p}(x)$ for a customer with
    balance $x=100$ EUR and for a customer with balance $x=250$ EUR.
  \item \pp{4} For the same two customers, compute the \emph{odds} of default
    $\hat{p}(x)/(1-\hat{p}(x))$ and the \emph{log-odds} (logit). Interpret the odds of the
    customer with $x=250$ in one sentence.
  \item \pp{3} Determine the balance $x^{*}$ at which the predicted default probability
    equals exactly $0.5$. Show your derivation.
  \item \pp{4} Interpret the estimated slope on the odds scale: by which \emph{factor} do the
    estimated odds of default change (i) per additional EUR of balance and (ii) per
    additional 100 EUR of balance? Why is the statement ``$\hat{\beta}_1=0.02$ means the
    default \emph{probability} rises by 0.02 per EUR'' wrong?
  \item \pp{3} Give two reasons why an ordinary linear regression of the 0/1 response $Y$
    on $x$ would be inappropriate here, even though it can be computed mechanically.
\end{enumerate}

\ifdefined\withsolutions
\begin{solutionbox}
\textbf{(a) Predicted probabilities.} The linear predictor (log-odds) is
$\hat{\eta}(x)=\hat{\beta}_0+\hat{\beta}_1x=-4+0.02x$.
\[
x=100:\ \hat{\eta}=-4+2=-2,\qquad
\hat{p}(100)=\frac{e^{-2}}{1+e^{-2}}=\frac{0.1353}{1.1353}=0.119 .
\]
\[
x=250:\ \hat{\eta}=-4+5=1,\qquad
\hat{p}(250)=\frac{e^{1}}{1+e^{1}}=\frac{2.7183}{3.7183}=0.731 .
\]

\medskip
\textbf{(b) Odds and log-odds.} Since
$\hat{p}/(1-\hat{p})=e^{\hat{\eta}}$ and $\log\big(\hat{p}/(1-\hat{p})\big)=\hat{\eta}$:
\begin{center}
\begin{tabular}{lcc}
\toprule
 & $x=100$ & $x=250$\\
\midrule
log-odds $\hat{\eta}(x)$ & $-2.000$ & $1.000$\\
odds $e^{\hat{\eta}(x)}$ & $e^{-2}=0.135$ & $e^{1}=2.718$\\
\bottomrule
\end{tabular}
\end{center}
Check via (a): $0.119/0.881=0.135$ and $0.731/0.269=2.718$. \checkmark\
\emph{Interpretation for $x=250$:} the odds are about $2.72$ to 1, i.e.\ default is
estimated to be about $2.7$ times as probable as non-default for this customer.

\medskip
\textbf{(c) Balance with $\hat{p}=0.5$.}
$\hat{p}(x^{*})=0.5$ $\iff$ odds $=1$ $\iff$ log-odds $=0$:
\[
-4+0.02\,x^{*}=0
\quad\Longrightarrow\quad
x^{*}=\frac{4}{0.02}=200\ \text{EUR}.
\]
Below 200 EUR the model predicts ``no default'' as more likely, above 200 EUR ``default''.

\medskip
\textbf{(d) Odds-scale interpretation.}
Per additional EUR the odds are multiplied by
$e^{\hat{\beta}_1}=e^{0.02}=1.020$, i.e.\ they increase by about $2.0\,\%$.
Per additional 100 EUR they are multiplied by
$e^{100\hat{\beta}_1}=e^{2}=7.389$ (an increase of about $639\,\%$).
The quoted statement is wrong because $\hat{\beta}_1$ is the change in the
\emph{log-odds} per EUR, not in the probability: the probability change
$\partial\hat{p}/\partial x=\hat{\beta}_1\hat{p}(x)\bigl(1-\hat{p}(x)\bigr)$ depends on $x$
--- the S-shaped curve is steepest near $\hat{p}=0.5$ and nearly flat in the tails, so there
is no constant ``probability effect per EUR''.

\medskip
\textbf{(e) Why not linear regression?}
(1) A straight line $\hat{\beta}_0+\hat{\beta}_1x$ is unbounded, so for small or large
balances it produces ``probabilities'' below 0 or above 1, which are not sensible
estimates of $\Pr(Y=1\mid x)$; the logistic function is guaranteed to stay in $[0,1]$.
(2) The 0/1 response violates the linear-model error assumptions: the errors are not
normal and their variance $p(x)(1-p(x))$ depends on $x$ (heteroscedasticity), so the usual
inference is unreliable. (Also, with more than two classes a numeric coding of $Y$ would
impose an artificial ordering and equal spacing between classes.)
\end{solutionbox}
\fi

% ============================================================
\problem{2: Discriminant analysis and naive Bayes --- concepts}{12}

\begin{enumerate}[label=(\alph*),leftmargin=7mm,itemsep=2.5mm]
  \item \pp{4} Linear discriminant analysis (LDA) and quadratic discriminant analysis (QDA)
    both model the class-conditional densities $f_k(x)$ of the predictors as multivariate
    normal. State precisely the assumption about the covariance matrices that distinguishes
    the two methods, and give the number of covariance \emph{matrices} each method estimates
    when there are $K$ classes.
  \item \pp{3} Using the bias--variance trade-off, explain in which situations LDA tends to
    outperform QDA on test data, and in which situations QDA tends to win. Mention the role
    of the sample size $n$ and of the true class covariances.
  \item \pp{3} State the key assumption of the \emph{naive Bayes} classifier about the
    predictors within each class, and explain why this assumption can improve test
    performance when the number of predictors $p$ is large relative to $n$ --- even though
    the assumption is rarely exactly true.
  \item \pp{2} Which of the three classifiers (LDA, QDA, naive Bayes) produce \emph{linear}
    decision boundaries in $x$, and which produce \emph{quadratic} or more general
    non-linear boundaries?
\end{enumerate}

\ifdefined\withsolutions
\begin{solutionbox}
\textbf{(a) Covariance assumptions.}
Both methods assume $X\mid Y=k\sim\mathcal{N}(\mu_k,\Sigma_k)$.
\emph{LDA} additionally assumes a covariance matrix that is \emph{shared} by all classes:
$\Sigma_1=\dots=\Sigma_K=\Sigma$; it therefore estimates \emph{one} common $p\times p$
covariance matrix. \emph{QDA} allows each class its \emph{own} covariance matrix
$\Sigma_k$ and estimates $K$ covariance matrices (one per class), i.e.\ roughly
$K\,p(p+1)/2$ covariance parameters instead of $p(p+1)/2$.

\medskip
\textbf{(b) When does which method win?}
QDA is the more flexible method (many more parameters): \emph{low bias, high variance};
LDA is more restrictive: \emph{higher bias if the $\Sigma_k$ truly differ, but much lower
variance}. Hence LDA tends to win when $n$ is small (or $p$ large), so that reducing
variance is crucial, and when the true class covariances are (approximately) equal ---
then LDA's assumption costs little bias. QDA tends to win when $n$ is large, so its extra
parameters can be estimated reliably, and when the class covariances clearly differ, so
LDA's shared-$\Sigma$ assumption creates substantial bias.

\medskip
\textbf{(c) Naive Bayes.}
Assumption: \emph{within each class the $p$ predictors are independent}, i.e.\ the joint
density factorises as $f_k(x)=f_{k1}(x_1)\cdot f_{k2}(x_2)\cdots f_{kp}(x_p)$.
Instead of a full $p$-dimensional joint density (which would require estimating
associations between all predictors, e.g.\ $p(p+1)/2$ covariance parameters per class),
only $p$ one-dimensional marginal densities per class must be estimated. This drastic
reduction in the number of parameters lowers the \emph{variance} of the classifier.
When $p$ is large relative to $n$, the variance saving usually outweighs the bias
introduced by (falsely) ignoring within-class dependence --- another instance of the
bias--variance trade-off.

\medskip
\textbf{(d) Shape of the decision boundaries.}
\emph{LDA:} linear boundaries (the discriminant functions are linear in $x$; the quadratic
term cancels because $\Sigma$ is shared). \emph{QDA:} quadratic boundaries (each class has
its own quadratic form). \emph{Naive Bayes:} generally non-linear boundaries of flexible
shape --- the log-odds are an \emph{additive} function $\sum_j g_j(x_j)$ of the individual
predictors (no interaction terms); only in special cases (e.g.\ Gaussian marginals with
class-shared variances) does it reduce to a linear rule.
\end{solutionbox}
\fi

% ============================================================
\problem{3: Confusion matrix and ROC curve}{15}

A logistic regression classifier predicts credit-card default on an evaluation set of
$n=1{,}000$ customers, using the default threshold: predict ``default'' if
$\hat{p}(\text{default}\mid x)>0.5$. The result is the following confusion matrix:

\begin{center}
\renewcommand{\arraystretch}{1.15}
\begin{tabular}{llcc|c}
 & & \multicolumn{2}{c|}{\textbf{True status}} & \\
 & & Default & No default & Total\\
\cmidrule{2-5}
\textbf{Predicted} & Default    & 60 & 45  & 105\\
                   & No default & 40 & 855 & 895\\
\cmidrule{2-5}
 & Total & 100 & 900 & 1{,}000\\
\end{tabular}
\end{center}

\begin{enumerate}[label=(\alph*),leftmargin=7mm,itemsep=2.5mm]
  \item \pp{7} Identify TP, FP, FN and TN, and compute (i) the overall accuracy and the
    overall error rate, (ii) the sensitivity (recall, true positive rate), (iii) the
    specificity, (iv) the precision, and (v) the false positive rate. Interpret sensitivity
    and precision in one sentence each in the context of this application.
  \item \pp{4} The bank now lowers the classification threshold from $0.5$ to $0.2$
    (predict ``default'' if $\hat{p}>0.2$). Explain --- without further computation --- how
    sensitivity and specificity will change, and give a business reason why the bank might
    deliberately accept this trade-off.
  \item \pp{4} The ROC curve plots which two quantities against each other, as what is
    varied? What does an area under the curve (AUC) of $0.5$ mean, and what does an AUC of
    $1$ mean?
\end{enumerate}

\ifdefined\withsolutions
\begin{solutionbox}
\textbf{(a) Metrics.} With ``default'' as the positive class:
$\text{TP}=60$, $\text{FP}=45$, $\text{FN}=40$, $\text{TN}=855$.
\begin{align*}
\text{(i) accuracy} &= \frac{\text{TP}+\text{TN}}{n}=\frac{60+855}{1000}=\frac{915}{1000}=0.915,
\qquad \text{error rate}=1-0.915=0.085,\\[0.5mm]
\text{(ii) sensitivity} &= \frac{\text{TP}}{\text{TP}+\text{FN}}=\frac{60}{100}=0.600,\\[0.5mm]
\text{(iii) specificity} &= \frac{\text{TN}}{\text{TN}+\text{FP}}=\frac{855}{900}=0.950,\\[0.5mm]
\text{(iv) precision} &= \frac{\text{TP}}{\text{TP}+\text{FP}}=\frac{60}{105}=0.571,\\[0.5mm]
\text{(v) FPR} &= \frac{\text{FP}}{\text{FP}+\text{TN}}=\frac{45}{900}=0.050=1-\text{specificity}.
\end{align*}
\emph{Sensitivity:} the classifier detects $60.0\,\%$ of the customers who actually
default --- $40\,\%$ of true defaulters are missed.
\emph{Precision:} of the customers flagged as defaulters, only $57.1\,\%$ actually
default; the rest are false alarms.

\medskip
\textbf{(b) Lowering the threshold to 0.2.}
More customers exceed the lower cut-off, so more are predicted ``default'': every true
defaulter with $0.2<\hat{p}\le 0.5$ is now caught, so \emph{sensitivity increases}
(FN $\downarrow$); at the same time more non-defaulters are flagged, so
\emph{specificity decreases} (FP $\uparrow$, FPR $\uparrow$). (The overall error rate will
typically increase, since non-defaulters are the large majority.)
\emph{Business reason:} the costs are asymmetric --- a missed defaulter (FN, a credit
loss) is far more expensive than a false alarm (FP, e.g.\ an unnecessary credit-line
review or customer call). Trading specificity for sensitivity can therefore minimise the
\emph{expected cost} even if the raw error rate rises.

\medskip
\textbf{(c) ROC and AUC.}
The ROC curve plots the \emph{true positive rate} (sensitivity) on the vertical axis
against the \emph{false positive rate} ($1-$specificity) on the horizontal axis, as the
classification \emph{threshold} is varied over all values from 1 down to 0.
$\text{AUC}=0.5$: performance equal to the diagonal, i.e.\ no better than random guessing
--- the score contains no usable ranking information. $\text{AUC}=1$: perfect
classifier --- some threshold separates the two classes without error
(sensitivity $=1$ and specificity $=1$ simultaneously). Equivalently, the AUC is the
probability that a randomly chosen defaulter receives a higher score $\hat{p}$ than a
randomly chosen non-defaulter.
\end{solutionbox}
\fi

% ============================================================
\problem{4: Resampling methods}{18}

\begin{enumerate}[label=(\alph*),leftmargin=7mm,itemsep=2.5mm]
  \item \pp{4} A data set has $n=200$ observations. To estimate the test error of a fixed
    model specification, state how many times the model must be \emph{fitted} for (i) the
    validation-set approach (one 50/50 split), (ii) 10-fold cross-validation, and (iii)
    leave-one-out cross-validation (LOOCV), with a one-sentence justification each.
  \item \pp{4} Compare LOOCV and 10-fold CV as estimators of the true test error with
    respect to \emph{bias} and \emph{variance}. Which method is preferred in practice, and
    why?
  \item \pp{5} A bootstrap sample of size $n$ is drawn with replacement from $n$
    observations. Show that the probability that a fixed observation $i$ is
    \emph{contained} in the bootstrap sample equals $1-(1-1/n)^n$. Evaluate this
    probability for $n=5$, and state the limiting value as $n\to\infty$
    (use $e^{-1}\approx 0.368$).
  \item \pp{5} To quantify the uncertainty of an estimated portfolio weight
    $\hat{\alpha}$, an analyst draws $B=3$ bootstrap samples (in practice $B\approx 1{,}000$;
    $B=3$ is used here only to keep the arithmetic manageable) and obtains
    \[
    \hat{\alpha}^{*1}=0.4,\qquad \hat{\alpha}^{*2}=0.5,\qquad \hat{\alpha}^{*3}=0.6 .
    \]
    Compute the bootstrap estimate of the standard error,
    \[
    \widehat{\operatorname{SE}}_B(\hat{\alpha})
      =\sqrt{\frac{1}{B-1}\sum_{r=1}^{B}\Bigl(\hat{\alpha}^{*r}-\bar{\alpha}^{*}\Bigr)^{2}},
    \qquad
    \bar{\alpha}^{*}=\frac{1}{B}\sum_{r=1}^{B}\hat{\alpha}^{*r},
    \]
    and explain in one sentence what this number estimates. Why is the bootstrap useful
    here although a formula for $\operatorname{SE}(\hat{\alpha})$ may be hard to derive?
\end{enumerate}

\ifdefined\withsolutions
\begin{solutionbox}
\textbf{(a) Number of model fits at $n=200$.}
(i) \emph{Validation set:} \textbf{1} fit --- the model is trained once on the training
half (100 observations) and evaluated on the held-out half.
(ii) \emph{10-fold CV:} \textbf{10} fits --- each of the 10 folds (20 observations) is
held out once, and the model is refitted on the remaining 180 observations each time.
(iii) \emph{LOOCV:} \textbf{200} fits --- each single observation is held out once and
the model is refitted on the other 199. (LOOCV $=$ $n$-fold CV; for least squares a
shortcut formula avoids the refits, but generically $n$ fits are needed.)

\medskip
\textbf{(b) Bias and variance of the CV estimators.}
\emph{Bias:} LOOCV trains on $n-1$ observations, almost the full sample, so it is nearly
unbiased for the true test error; 10-fold CV trains on $0.9\,n$ and therefore slightly
\emph{overestimates} the test error (a bit more bias). \emph{Variance:} LOOCV has the
\emph{higher} variance: the $n$ fitted models are trained on almost identical data sets,
so their errors are strongly positively correlated, and the mean of highly correlated
quantities has high variance; the 10 folds of 10-fold CV overlap less, giving less
correlated fits and a lower-variance estimate. \emph{In practice} $k=5$ or $k=10$ is
preferred: a good bias--variance compromise and much cheaper computationally.

\medskip
\textbf{(c) Bootstrap inclusion probability.}
Each of the $n$ draws is made with replacement, so each draw selects observation $i$ with
probability $1/n$ and misses it with probability $1-1/n$. The draws are independent, so
\[
\Pr(i\ \text{not in sample})=\Bigl(1-\tfrac{1}{n}\Bigr)^{n}
\quad\Longrightarrow\quad
\Pr(i\ \text{in sample})=1-\Bigl(1-\tfrac{1}{n}\Bigr)^{n}.
\]
For $n=5$: $1-(4/5)^5=1-0.32768=0.672$.
As $n\to\infty$: $(1-1/n)^n\to e^{-1}$, so the probability tends to
$1-e^{-1}\approx 1-0.368=0.632$ --- a bootstrap sample contains about $63\,\%$ of the
distinct original observations, even for large $n$.

\medskip
\textbf{(d) Bootstrap standard error.}
$\bar{\alpha}^{*}=(0.4+0.5+0.6)/3=0.5$.
\[
\widehat{\operatorname{SE}}_B(\hat{\alpha})
=\sqrt{\frac{(0.4-0.5)^2+(0.5-0.5)^2+(0.6-0.5)^2}{3-1}}
=\sqrt{\frac{0.01+0+0.01}{2}}
=\sqrt{0.01}=0.100 .
\]
This estimates the standard deviation of $\hat{\alpha}$ over repeated samples from the
population, i.e.\ how much the estimated weight would typically vary if the data could be
collected anew. The bootstrap is useful because it replaces analytic derivations by
resampling from the observed data: it needs no formula for
$\operatorname{SE}(\hat{\alpha})$ and no new data, and it works for almost any statistic
--- including complicated ones for which no closed-form standard error exists.
\end{solutionbox}
\fi

% ============================================================
\problem{5: Model selection with $C_p$ and BIC}{12}

Best subset selection on a data set with $n=100$ observations produced, for each model
size $d$ (number of predictors), the best model and its residual sum of squares:

\begin{center}
\begin{tabular}{lcccc}
\toprule
Number of predictors $d$ & 1 & 2 & 3 & 4\\
\midrule
$\operatorname{RSS}_d$ & 500 & 444 & 428 & 424\\
\bottomrule
\end{tabular}
\end{center}

The noise variance is estimated from the full model as $\hat{\sigma}^2=4$.
Use the criteria
\[
C_p=\frac{1}{n}\bigl(\operatorname{RSS}+2\,d\,\hat{\sigma}^2\bigr),
\qquad
\operatorname{BIC}=\frac{1}{n}\bigl(\operatorname{RSS}+\log(n)\,d\,\hat{\sigma}^2\bigr),
\qquad \log(100)\approx 4.605 .
\]

\begin{enumerate}[label=(\alph*),leftmargin=7mm,itemsep=2.5mm]
  \item \pp{4} Compute $C_p$ for all four models (3 decimals).
  \item \pp{4} Compute BIC for all four models (3 decimals).
  \item \pp{2} Which model size does each criterion select? Explain the reason for the
    disagreement by comparing the two penalty terms.
  \item \pp{2} Why can the training RSS (or equivalently the training $R^2$) alone never be
    used to select $d$? What does RSS necessarily do as $d$ grows?
\end{enumerate}

\ifdefined\withsolutions
\begin{solutionbox}
\textbf{(a) $C_p$.} Penalty per predictor: $2\hat{\sigma}^2=8$; $C_p=(\operatorname{RSS}_d+8d)/100$:
\begin{center}
\begin{tabular}{lcccc}
\toprule
$d$ & 1 & 2 & 3 & 4\\
\midrule
$\operatorname{RSS}_d+8d$ & 508 & 460 & 452 & 456\\
$C_p$ & 5.080 & 4.600 & \textbf{4.520} & 4.560\\
\bottomrule
\end{tabular}
\end{center}

\medskip
\textbf{(b) BIC.} Penalty per predictor: $\log(100)\,\hat{\sigma}^2=4.605\cdot 4=18.421$;
$\operatorname{BIC}=(\operatorname{RSS}_d+18.421\,d)/100$:
\begin{center}
\begin{tabular}{lcccc}
\toprule
$d$ & 1 & 2 & 3 & 4\\
\midrule
$\operatorname{RSS}_d+18.421\,d$ & 518.421 & 480.841 & 483.262 & 497.683\\
BIC & 5.184 & \textbf{4.808} & 4.833 & 4.977\\
\bottomrule
\end{tabular}
\end{center}

\medskip
\textbf{(c) Selected models.}
$C_p$ is minimised at $d=3$ (4.520); BIC is minimised at $d=2$ (4.808).
Reason: both criteria are RSS plus a penalty for model size, but BIC charges
$\log(n)\hat{\sigma}^2=18.421$ per predictor whereas $C_p$ charges only
$2\hat{\sigma}^2=8$; since $\log(100)=4.605>2$, BIC penalises additional predictors much
more heavily and therefore tends to select \emph{smaller} models. Here the RSS drop from
$d=2$ to $d=3$ (16) exceeds the $C_p$ penalty (8) but not the BIC penalty (18.421).

\medskip
\textbf{(d) Why RSS alone cannot choose $d$.}
For nested model searches, adding a predictor can never increase the training RSS: RSS
decreases \emph{monotonically} in $d$ (and training $R^2$ increases monotonically).
RSS-based selection would therefore always pick the largest model, $d=4$, regardless of
whether the extra predictors carry real signal --- training error is an ever more
optimistic estimate of test error as flexibility grows. Hence one must either penalise
model size ($C_p$, AIC, BIC, adjusted $R^2$) or estimate the test error directly
(validation set, cross-validation).
\end{solutionbox}
\fi

% ============================================================
\problem{6: Ridge regression and the lasso}{15}

Consider a linear model with response $y$ and $p$ standardised predictors
$X_1,\dots,X_p$. Part (e) also draws on the bias--variance ideas of Chapter~2 and the
least-squares fit of Chapter~3.

\begin{enumerate}[label=(\alph*),leftmargin=7mm,itemsep=2.5mm]
  \item \pp{4} Write down the optimisation objectives of \emph{ridge regression} and of
    the \emph{lasso} (RSS plus the respective penalty, with tuning parameter
    $\lambda\ge 0$). Which coefficient is conventionally \emph{not} penalised?
  \item \pp{3} Describe the behaviour of the two estimators in the limits
    $\lambda\to 0$ and $\lambda\to\infty$. Which method from Chapter~3 is recovered as
    $\lambda\to 0$?
  \item \pp{3} Explain --- in words, using the constraint-region geometry for $p=2$ --- why
    the lasso can set coefficients \emph{exactly} to zero while ridge regression
    typically does not.
  \item \pp{2} Why should the predictors be \emph{standardised} before ridge or lasso is
    applied, when this is irrelevant for ordinary least squares?
  \item \pp{3} A ridge regression was fitted for two values of the tuning parameter; the
    table shows the standardised coefficient estimates. A 10-fold cross-validation over a
    grid of $\lambda$ values gave the CV errors on the right.
    \begin{center}
    \begin{tabular}{lrr}
    \toprule
     & $\lambda=0.1$ & $\lambda=100$\\
    \midrule
    $\hat{\beta}_1$ & $0.92$ & $0.31$\\
    $\hat{\beta}_2$ & $0.85$ & $0.29$\\
    $\hat{\beta}_3$ & $-0.41$ & $-0.12$\\
    $\hat{\beta}_4$ & $0.20$ & $0.06$\\
    \bottomrule
    \end{tabular}
    \hspace{12mm}
    \begin{tabular}{lcccc}
    \toprule
    $\lambda$ & 0.1 & 1 & 10 & 100\\
    \midrule
    CV error & 12.8 & 11.9 & 10.6 & 13.4\\
    \bottomrule
    \end{tabular}
    \end{center}
    Which of the two fits shown in the coefficient table has the lower bias and which the
    lower variance, and why? Which $\lambda$ from the CV table should be chosen for the
    final model?
\end{enumerate}

\ifdefined\withsolutions
\begin{solutionbox}
\textbf{(a) Objectives.} Ridge regression minimises
\[
\sum_{i=1}^{n}\Bigl(y_i-\beta_0-\sum_{j=1}^{p}\beta_j x_{ij}\Bigr)^{2}
+\lambda\sum_{j=1}^{p}\beta_j^{2}
=\operatorname{RSS}+\lambda\|\beta\|_2^2 ,
\]
the lasso minimises
\[
\sum_{i=1}^{n}\Bigl(y_i-\beta_0-\sum_{j=1}^{p}\beta_j x_{ij}\Bigr)^{2}
+\lambda\sum_{j=1}^{p}|\beta_j|
=\operatorname{RSS}+\lambda\|\beta\|_1 .
\]
The intercept $\beta_0$ is \emph{not} penalised (shrinking the overall level of $y$ would
make no sense).

\medskip
\textbf{(b) Limiting behaviour.}
$\lambda\to 0$: the penalty vanishes and both estimators approach the \emph{ordinary
least-squares} solution of Chapter~3.
$\lambda\to\infty$: the penalty dominates and all penalised coefficients are forced to
zero --- the null model containing only the intercept. Ridge shrinks the coefficients
\emph{towards} zero (exactly zero only in the limit), whereas the lasso sets more and more
coefficients \emph{exactly} to zero as $\lambda$ grows. In between, $\lambda$ traces out a
whole path of fits whose flexibility decreases as $\lambda$ increases.

\medskip
\textbf{(c) Geometry of the penalties ($p=2$).}
The penalised problems are equivalent to minimising the RSS subject to a budget
constraint: $|\beta_1|+|\beta_2|\le s$ (lasso) or $\beta_1^2+\beta_2^2\le s$ (ridge).
The lasso region is a \emph{diamond} with corners on the coordinate axes; the ridge region
is a \emph{disk} with no corners. The RSS is minimised over the constraint region where
its elliptical contours around the OLS solution first touch the region. Ellipses often hit
the diamond exactly at a \emph{corner} --- and at a corner one coefficient is exactly zero
(sparsity, variable selection). The circle has no corners, so the touching point generically
has both coordinates non-zero: ridge shrinks but keeps all $p$ predictors.

\medskip
\textbf{(d) Standardisation.}
OLS is scale-equivariant: rescaling $X_j$ (e.g.\ EUR $\to$ thousand EUR) simply rescales
$\hat{\beta}_j$, leaving fits unchanged. The penalties, however, treat all coefficients
symmetrically: $\lambda\sum\beta_j^2$ resp.\ $\lambda\sum|\beta_j|$ depend on the
\emph{magnitudes} of the coefficients and hence on the units of the predictors --- a
predictor measured in small units has a large coefficient and would be penalised
(shrunk) more strongly for purely cosmetic reasons. Standardising each predictor to
standard deviation one makes the penalty act comparably on all predictors and the solution
independent of the original units.

\medskip
\textbf{(e) Bias--variance and choice of $\lambda$.}
$\lambda=0.1$: almost no shrinkage --- the fit is close to OLS, i.e.\ more flexible:
\emph{lower bias, higher variance}. $\lambda=100$: strong shrinkage (coefficients roughly
one third of their former size, e.g.\ $0.92\to 0.31$): a more rigid fit with
\emph{higher bias but lower variance} --- the bias--variance trade-off of Chapter~2
governed by a single knob $\lambda$.
From the CV table, the estimated test error is minimised at $\lambda=10$
(CV error $10.6<11.9<12.8<13.4$), so \emph{$\lambda=10$} should be chosen --- an
intermediate amount of shrinkage; neither of the two fits shown in the coefficient table
is optimal.
\end{solutionbox}
\fi

\vspace{6mm}
\begin{center}
\rule{0.35\textwidth}{0.4pt}\\[1mm]
\textit{End of the examination --- good luck!}
\end{center}

\end{document}
"
% ============================================================
\documentclass[11pt,a4paper]{article}
\usepackage[utf8]{inputenc}
\usepackage[T1]{fontenc}
\usepackage[english]{babel}
\usepackage[margin=2.2cm]{geometry}
\usepackage{amsmath,amssymb}
\usepackage{booktabs,array}
\usepackage{enumitem}
\usepackage{xcolor}
\usepackage{tcolorbox}
\tcbuselibrary{skins,breakable}

\setlength{\parindent}{0pt}

% Teal solution box (only shown when \withsolutions is defined)
\newtcolorbox{solutionbox}{enhanced,breakable,colback=teal!5,colframe=teal!55!black,
  boxrule=0pt,leftrule=2.5pt,arc=2pt,fonttitle=\bfseries,title=Solution,
  left=2mm,right=2mm,top=1mm,bottom=1.5mm}

\newcommand{\problem}[2]{\par\vspace{7mm}\noindent{\large\textbf{Problem #1}}%
  \hfill{\large\textbf{(#2 points)}}\par\nopagebreak\vspace{0.6mm}%
  \noindent\rule{\textwidth}{0.6pt}\par\nopagebreak\vspace{2.5mm}}
\newcommand{\pp}[1]{\textbf{[#1~P]}\enspace}

\begin{document}

% ------------------------------------------------------------
% Header block
% ------------------------------------------------------------
\begin{center}
  {\Large\textbf{HSBI --- Bielefeld University of Applied Sciences and Arts}}\\[1.2mm]
  {\large Quantitative Research Methods}\\[1mm]
  {\normalsize Summer Semester 2026 \quad$\cdot$\quad Examiner: Prof.\ Dr.\ Christoph Weisser}\\[4.5mm]
  {\LARGE\textbf{Mock Examination 2 (Lectures 5--8)}}\\[1.2mm]
  {\normalsize ISLP Chapters 4--6: Classification $\cdot$ Resampling Methods $\cdot$ Model Selection and Regularisation}%
  \ifdefined\withsolutions\\[2.2mm]{\large\textcolor{teal!55!black}{\textbf{--- Version with detailed solutions ---}}}\fi
\end{center}

\vspace{2.5mm}
\begin{center}
\begin{tabular}{@{}ll@{}}
  \textbf{Duration:} & \textbf{90 minutes}\\
  \textbf{Total credit:} & \textbf{90 points}\\
  \textbf{Allowed aids:} & non-programmable calculator; one handwritten A4 sheet of notes (both sides)\\
\end{tabular}
\end{center}

\vspace{3mm}
Name: \rule[-1mm]{0.38\textwidth}{0.4pt}\hfill Matriculation no.: \rule[-1mm]{0.24\textwidth}{0.4pt}

\vspace{4.5mm}
\begin{center}
\renewcommand{\arraystretch}{1.35}
\begin{tabular}{|l|c|c|c|c|c|c|c|}
\hline
\textbf{Problem} & \textbf{P1} & \textbf{P2} & \textbf{P3} & \textbf{P4} & \textbf{P5} & \textbf{P6} & \textbf{Total}\\\hline
Maximum points & 18 & 12 & 15 & 18 & 12 & 15 & \textbf{90}\\\hline
Achieved points & \hspace{9mm} & \hspace{9mm} & \hspace{9mm} & \hspace{9mm} & \hspace{9mm} & \hspace{9mm} & \hspace{11mm}\\\hline
\end{tabular}
\end{center}

\vspace{2.5mm}
\textbf{Instructions}
\begin{itemize}[nosep,leftmargin=5.5mm]
  \item Justify all answers. Numerical results without a derivation or a brief reasoning receive no credit.
  \item Round final numerical results to 3 decimal places unless stated otherwise. Small deviations caused by carrying rounded intermediate results are accepted.
  \item The number of points per problem roughly equals the recommended working time in minutes.
\end{itemize}
\vspace{1mm}\noindent\rule{\textwidth}{0.6pt}

% ============================================================
\problem{1: Logistic regression by hand}{18}

A bank models the probability that a credit-card customer \emph{defaults} ($Y=1$) as a
function of the monthly card balance $x$ (in EUR; balances in this portfolio range from
0 to 400 EUR). A logistic regression fitted by maximum likelihood gave
\[
\hat{p}(x)=\Pr(Y=1\mid X=x)
  =\frac{e^{\hat{\beta}_0+\hat{\beta}_1 x}}{1+e^{\hat{\beta}_0+\hat{\beta}_1 x}},
\qquad
\hat{\beta}_0=-4,\quad \hat{\beta}_1=0.02 .
\]

\begin{enumerate}[label=(\alph*),leftmargin=7mm,itemsep=2.5mm]
  \item \pp{4} Compute the predicted default probability $\hat{p}(x)$ for a customer with
    balance $x=100$ EUR and for a customer with balance $x=250$ EUR.
  \item \pp{4} For the same two customers, compute the \emph{odds} of default
    $\hat{p}(x)/(1-\hat{p}(x))$ and the \emph{log-odds} (logit). Interpret the odds of the
    customer with $x=250$ in one sentence.
  \item \pp{3} Determine the balance $x^{*}$ at which the predicted default probability
    equals exactly $0.5$. Show your derivation.
  \item \pp{4} Interpret the estimated slope on the odds scale: by which \emph{factor} do the
    estimated odds of default change (i) per additional EUR of balance and (ii) per
    additional 100 EUR of balance? Why is the statement ``$\hat{\beta}_1=0.02$ means the
    default \emph{probability} rises by 0.02 per EUR'' wrong?
  \item \pp{3} Give two reasons why an ordinary linear regression of the 0/1 response $Y$
    on $x$ would be inappropriate here, even though it can be computed mechanically.
\end{enumerate}

\ifdefined\withsolutions
\begin{solutionbox}
\textbf{(a) Predicted probabilities.} The linear predictor (log-odds) is
$\hat{\eta}(x)=\hat{\beta}_0+\hat{\beta}_1x=-4+0.02x$.
\[
x=100:\ \hat{\eta}=-4+2=-2,\qquad
\hat{p}(100)=\frac{e^{-2}}{1+e^{-2}}=\frac{0.1353}{1.1353}=0.119 .
\]
\[
x=250:\ \hat{\eta}=-4+5=1,\qquad
\hat{p}(250)=\frac{e^{1}}{1+e^{1}}=\frac{2.7183}{3.7183}=0.731 .
\]

\medskip
\textbf{(b) Odds and log-odds.} Since
$\hat{p}/(1-\hat{p})=e^{\hat{\eta}}$ and $\log\big(\hat{p}/(1-\hat{p})\big)=\hat{\eta}$:
\begin{center}
\begin{tabular}{lcc}
\toprule
 & $x=100$ & $x=250$\\
\midrule
log-odds $\hat{\eta}(x)$ & $-2.000$ & $1.000$\\
odds $e^{\hat{\eta}(x)}$ & $e^{-2}=0.135$ & $e^{1}=2.718$\\
\bottomrule
\end{tabular}
\end{center}
Check via (a): $0.119/0.881=0.135$ and $0.731/0.269=2.718$. \checkmark\
\emph{Interpretation for $x=250$:} the odds are about $2.72$ to 1, i.e.\ default is
estimated to be about $2.7$ times as probable as non-default for this customer.

\medskip
\textbf{(c) Balance with $\hat{p}=0.5$.}
$\hat{p}(x^{*})=0.5$ $\iff$ odds $=1$ $\iff$ log-odds $=0$:
\[
-4+0.02\,x^{*}=0
\quad\Longrightarrow\quad
x^{*}=\frac{4}{0.02}=200\ \text{EUR}.
\]
Below 200 EUR the model predicts ``no default'' as more likely, above 200 EUR ``default''.

\medskip
\textbf{(d) Odds-scale interpretation.}
Per additional EUR the odds are multiplied by
$e^{\hat{\beta}_1}=e^{0.02}=1.020$, i.e.\ they increase by about $2.0\,\%$.
Per additional 100 EUR they are multiplied by
$e^{100\hat{\beta}_1}=e^{2}=7.389$ (an increase of about $639\,\%$).
The quoted statement is wrong because $\hat{\beta}_1$ is the change in the
\emph{log-odds} per EUR, not in the probability: the probability change
$\partial\hat{p}/\partial x=\hat{\beta}_1\hat{p}(x)\bigl(1-\hat{p}(x)\bigr)$ depends on $x$
--- the S-shaped curve is steepest near $\hat{p}=0.5$ and nearly flat in the tails, so there
is no constant ``probability effect per EUR''.

\medskip
\textbf{(e) Why not linear regression?}
(1) A straight line $\hat{\beta}_0+\hat{\beta}_1x$ is unbounded, so for small or large
balances it produces ``probabilities'' below 0 or above 1, which are not sensible
estimates of $\Pr(Y=1\mid x)$; the logistic function is guaranteed to stay in $[0,1]$.
(2) The 0/1 response violates the linear-model error assumptions: the errors are not
normal and their variance $p(x)(1-p(x))$ depends on $x$ (heteroscedasticity), so the usual
inference is unreliable. (Also, with more than two classes a numeric coding of $Y$ would
impose an artificial ordering and equal spacing between classes.)
\end{solutionbox}
\fi

% ============================================================
\problem{2: Discriminant analysis and naive Bayes --- concepts}{12}

\begin{enumerate}[label=(\alph*),leftmargin=7mm,itemsep=2.5mm]
  \item \pp{4} Linear discriminant analysis (LDA) and quadratic discriminant analysis (QDA)
    both model the class-conditional densities $f_k(x)$ of the predictors as multivariate
    normal. State precisely the assumption about the covariance matrices that distinguishes
    the two methods, and give the number of covariance \emph{matrices} each method estimates
    when there are $K$ classes.
  \item \pp{3} Using the bias--variance trade-off, explain in which situations LDA tends to
    outperform QDA on test data, and in which situations QDA tends to win. Mention the role
    of the sample size $n$ and of the true class covariances.
  \item \pp{3} State the key assumption of the \emph{naive Bayes} classifier about the
    predictors within each class, and explain why this assumption can improve test
    performance when the number of predictors $p$ is large relative to $n$ --- even though
    the assumption is rarely exactly true.
  \item \pp{2} Which of the three classifiers (LDA, QDA, naive Bayes) produce \emph{linear}
    decision boundaries in $x$, and which produce \emph{quadratic} or more general
    non-linear boundaries?
\end{enumerate}

\ifdefined\withsolutions
\begin{solutionbox}
\textbf{(a) Covariance assumptions.}
Both methods assume $X\mid Y=k\sim\mathcal{N}(\mu_k,\Sigma_k)$.
\emph{LDA} additionally assumes a covariance matrix that is \emph{shared} by all classes:
$\Sigma_1=\dots=\Sigma_K=\Sigma$; it therefore estimates \emph{one} common $p\times p$
covariance matrix. \emph{QDA} allows each class its \emph{own} covariance matrix
$\Sigma_k$ and estimates $K$ covariance matrices (one per class), i.e.\ roughly
$K\,p(p+1)/2$ covariance parameters instead of $p(p+1)/2$.

\medskip
\textbf{(b) When does which method win?}
QDA is the more flexible method (many more parameters): \emph{low bias, high variance};
LDA is more restrictive: \emph{higher bias if the $\Sigma_k$ truly differ, but much lower
variance}. Hence LDA tends to win when $n$ is small (or $p$ large), so that reducing
variance is crucial, and when the true class covariances are (approximately) equal ---
then LDA's assumption costs little bias. QDA tends to win when $n$ is large, so its extra
parameters can be estimated reliably, and when the class covariances clearly differ, so
LDA's shared-$\Sigma$ assumption creates substantial bias.

\medskip
\textbf{(c) Naive Bayes.}
Assumption: \emph{within each class the $p$ predictors are independent}, i.e.\ the joint
density factorises as $f_k(x)=f_{k1}(x_1)\cdot f_{k2}(x_2)\cdots f_{kp}(x_p)$.
Instead of a full $p$-dimensional joint density (which would require estimating
associations between all predictors, e.g.\ $p(p+1)/2$ covariance parameters per class),
only $p$ one-dimensional marginal densities per class must be estimated. This drastic
reduction in the number of parameters lowers the \emph{variance} of the classifier.
When $p$ is large relative to $n$, the variance saving usually outweighs the bias
introduced by (falsely) ignoring within-class dependence --- another instance of the
bias--variance trade-off.

\medskip
\textbf{(d) Shape of the decision boundaries.}
\emph{LDA:} linear boundaries (the discriminant functions are linear in $x$; the quadratic
term cancels because $\Sigma$ is shared). \emph{QDA:} quadratic boundaries (each class has
its own quadratic form). \emph{Naive Bayes:} generally non-linear boundaries of flexible
shape --- the log-odds are an \emph{additive} function $\sum_j g_j(x_j)$ of the individual
predictors (no interaction terms); only in special cases (e.g.\ Gaussian marginals with
class-shared variances) does it reduce to a linear rule.
\end{solutionbox}
\fi

% ============================================================
\problem{3: Confusion matrix and ROC curve}{15}

A logistic regression classifier predicts credit-card default on an evaluation set of
$n=1{,}000$ customers, using the default threshold: predict ``default'' if
$\hat{p}(\text{default}\mid x)>0.5$. The result is the following confusion matrix:

\begin{center}
\renewcommand{\arraystretch}{1.15}
\begin{tabular}{llcc|c}
 & & \multicolumn{2}{c|}{\textbf{True status}} & \\
 & & Default & No default & Total\\
\cmidrule{2-5}
\textbf{Predicted} & Default    & 60 & 45  & 105\\
                   & No default & 40 & 855 & 895\\
\cmidrule{2-5}
 & Total & 100 & 900 & 1{,}000\\
\end{tabular}
\end{center}

\begin{enumerate}[label=(\alph*),leftmargin=7mm,itemsep=2.5mm]
  \item \pp{7} Identify TP, FP, FN and TN, and compute (i) the overall accuracy and the
    overall error rate, (ii) the sensitivity (recall, true positive rate), (iii) the
    specificity, (iv) the precision, and (v) the false positive rate. Interpret sensitivity
    and precision in one sentence each in the context of this application.
  \item \pp{4} The bank now lowers the classification threshold from $0.5$ to $0.2$
    (predict ``default'' if $\hat{p}>0.2$). Explain --- without further computation --- how
    sensitivity and specificity will change, and give a business reason why the bank might
    deliberately accept this trade-off.
  \item \pp{4} The ROC curve plots which two quantities against each other, as what is
    varied? What does an area under the curve (AUC) of $0.5$ mean, and what does an AUC of
    $1$ mean?
\end{enumerate}

\ifdefined\withsolutions
\begin{solutionbox}
\textbf{(a) Metrics.} With ``default'' as the positive class:
$\text{TP}=60$, $\text{FP}=45$, $\text{FN}=40$, $\text{TN}=855$.
\begin{align*}
\text{(i) accuracy} &= \frac{\text{TP}+\text{TN}}{n}=\frac{60+855}{1000}=\frac{915}{1000}=0.915,
\qquad \text{error rate}=1-0.915=0.085,\\[0.5mm]
\text{(ii) sensitivity} &= \frac{\text{TP}}{\text{TP}+\text{FN}}=\frac{60}{100}=0.600,\\[0.5mm]
\text{(iii) specificity} &= \frac{\text{TN}}{\text{TN}+\text{FP}}=\frac{855}{900}=0.950,\\[0.5mm]
\text{(iv) precision} &= \frac{\text{TP}}{\text{TP}+\text{FP}}=\frac{60}{105}=0.571,\\[0.5mm]
\text{(v) FPR} &= \frac{\text{FP}}{\text{FP}+\text{TN}}=\frac{45}{900}=0.050=1-\text{specificity}.
\end{align*}
\emph{Sensitivity:} the classifier detects $60.0\,\%$ of the customers who actually
default --- $40\,\%$ of true defaulters are missed.
\emph{Precision:} of the customers flagged as defaulters, only $57.1\,\%$ actually
default; the rest are false alarms.

\medskip
\textbf{(b) Lowering the threshold to 0.2.}
More customers exceed the lower cut-off, so more are predicted ``default'': every true
defaulter with $0.2<\hat{p}\le 0.5$ is now caught, so \emph{sensitivity increases}
(FN $\downarrow$); at the same time more non-defaulters are flagged, so
\emph{specificity decreases} (FP $\uparrow$, FPR $\uparrow$). (The overall error rate will
typically increase, since non-defaulters are the large majority.)
\emph{Business reason:} the costs are asymmetric --- a missed defaulter (FN, a credit
loss) is far more expensive than a false alarm (FP, e.g.\ an unnecessary credit-line
review or customer call). Trading specificity for sensitivity can therefore minimise the
\emph{expected cost} even if the raw error rate rises.

\medskip
\textbf{(c) ROC and AUC.}
The ROC curve plots the \emph{true positive rate} (sensitivity) on the vertical axis
against the \emph{false positive rate} ($1-$specificity) on the horizontal axis, as the
classification \emph{threshold} is varied over all values from 1 down to 0.
$\text{AUC}=0.5$: performance equal to the diagonal, i.e.\ no better than random guessing
--- the score contains no usable ranking information. $\text{AUC}=1$: perfect
classifier --- some threshold separates the two classes without error
(sensitivity $=1$ and specificity $=1$ simultaneously). Equivalently, the AUC is the
probability that a randomly chosen defaulter receives a higher score $\hat{p}$ than a
randomly chosen non-defaulter.
\end{solutionbox}
\fi

% ============================================================
\problem{4: Resampling methods}{18}

\begin{enumerate}[label=(\alph*),leftmargin=7mm,itemsep=2.5mm]
  \item \pp{4} A data set has $n=200$ observations. To estimate the test error of a fixed
    model specification, state how many times the model must be \emph{fitted} for (i) the
    validation-set approach (one 50/50 split), (ii) 10-fold cross-validation, and (iii)
    leave-one-out cross-validation (LOOCV), with a one-sentence justification each.
  \item \pp{4} Compare LOOCV and 10-fold CV as estimators of the true test error with
    respect to \emph{bias} and \emph{variance}. Which method is preferred in practice, and
    why?
  \item \pp{5} A bootstrap sample of size $n$ is drawn with replacement from $n$
    observations. Show that the probability that a fixed observation $i$ is
    \emph{contained} in the bootstrap sample equals $1-(1-1/n)^n$. Evaluate this
    probability for $n=5$, and state the limiting value as $n\to\infty$
    (use $e^{-1}\approx 0.368$).
  \item \pp{5} To quantify the uncertainty of an estimated portfolio weight
    $\hat{\alpha}$, an analyst draws $B=3$ bootstrap samples (in practice $B\approx 1{,}000$;
    $B=3$ is used here only to keep the arithmetic manageable) and obtains
    \[
    \hat{\alpha}^{*1}=0.4,\qquad \hat{\alpha}^{*2}=0.5,\qquad \hat{\alpha}^{*3}=0.6 .
    \]
    Compute the bootstrap estimate of the standard error,
    \[
    \widehat{\operatorname{SE}}_B(\hat{\alpha})
      =\sqrt{\frac{1}{B-1}\sum_{r=1}^{B}\Bigl(\hat{\alpha}^{*r}-\bar{\alpha}^{*}\Bigr)^{2}},
    \qquad
    \bar{\alpha}^{*}=\frac{1}{B}\sum_{r=1}^{B}\hat{\alpha}^{*r},
    \]
    and explain in one sentence what this number estimates. Why is the bootstrap useful
    here although a formula for $\operatorname{SE}(\hat{\alpha})$ may be hard to derive?
\end{enumerate}

\ifdefined\withsolutions
\begin{solutionbox}
\textbf{(a) Number of model fits at $n=200$.}
(i) \emph{Validation set:} \textbf{1} fit --- the model is trained once on the training
half (100 observations) and evaluated on the held-out half.
(ii) \emph{10-fold CV:} \textbf{10} fits --- each of the 10 folds (20 observations) is
held out once, and the model is refitted on the remaining 180 observations each time.
(iii) \emph{LOOCV:} \textbf{200} fits --- each single observation is held out once and
the model is refitted on the other 199. (LOOCV $=$ $n$-fold CV; for least squares a
shortcut formula avoids the refits, but generically $n$ fits are needed.)

\medskip
\textbf{(b) Bias and variance of the CV estimators.}
\emph{Bias:} LOOCV trains on $n-1$ observations, almost the full sample, so it is nearly
unbiased for the true test error; 10-fold CV trains on $0.9\,n$ and therefore slightly
\emph{overestimates} the test error (a bit more bias). \emph{Variance:} LOOCV has the
\emph{higher} variance: the $n$ fitted models are trained on almost identical data sets,
so their errors are strongly positively correlated, and the mean of highly correlated
quantities has high variance; the 10 folds of 10-fold CV overlap less, giving less
correlated fits and a lower-variance estimate. \emph{In practice} $k=5$ or $k=10$ is
preferred: a good bias--variance compromise and much cheaper computationally.

\medskip
\textbf{(c) Bootstrap inclusion probability.}
Each of the $n$ draws is made with replacement, so each draw selects observation $i$ with
probability $1/n$ and misses it with probability $1-1/n$. The draws are independent, so
\[
\Pr(i\ \text{not in sample})=\Bigl(1-\tfrac{1}{n}\Bigr)^{n}
\quad\Longrightarrow\quad
\Pr(i\ \text{in sample})=1-\Bigl(1-\tfrac{1}{n}\Bigr)^{n}.
\]
For $n=5$: $1-(4/5)^5=1-0.32768=0.672$.
As $n\to\infty$: $(1-1/n)^n\to e^{-1}$, so the probability tends to
$1-e^{-1}\approx 1-0.368=0.632$ --- a bootstrap sample contains about $63\,\%$ of the
distinct original observations, even for large $n$.

\medskip
\textbf{(d) Bootstrap standard error.}
$\bar{\alpha}^{*}=(0.4+0.5+0.6)/3=0.5$.
\[
\widehat{\operatorname{SE}}_B(\hat{\alpha})
=\sqrt{\frac{(0.4-0.5)^2+(0.5-0.5)^2+(0.6-0.5)^2}{3-1}}
=\sqrt{\frac{0.01+0+0.01}{2}}
=\sqrt{0.01}=0.100 .
\]
This estimates the standard deviation of $\hat{\alpha}$ over repeated samples from the
population, i.e.\ how much the estimated weight would typically vary if the data could be
collected anew. The bootstrap is useful because it replaces analytic derivations by
resampling from the observed data: it needs no formula for
$\operatorname{SE}(\hat{\alpha})$ and no new data, and it works for almost any statistic
--- including complicated ones for which no closed-form standard error exists.
\end{solutionbox}
\fi

% ============================================================
\problem{5: Model selection with $C_p$ and BIC}{12}

Best subset selection on a data set with $n=100$ observations produced, for each model
size $d$ (number of predictors), the best model and its residual sum of squares:

\begin{center}
\begin{tabular}{lcccc}
\toprule
Number of predictors $d$ & 1 & 2 & 3 & 4\\
\midrule
$\operatorname{RSS}_d$ & 500 & 444 & 428 & 424\\
\bottomrule
\end{tabular}
\end{center}

The noise variance is estimated from the full model as $\hat{\sigma}^2=4$.
Use the criteria
\[
C_p=\frac{1}{n}\bigl(\operatorname{RSS}+2\,d\,\hat{\sigma}^2\bigr),
\qquad
\operatorname{BIC}=\frac{1}{n}\bigl(\operatorname{RSS}+\log(n)\,d\,\hat{\sigma}^2\bigr),
\qquad \log(100)\approx 4.605 .
\]

\begin{enumerate}[label=(\alph*),leftmargin=7mm,itemsep=2.5mm]
  \item \pp{4} Compute $C_p$ for all four models (3 decimals).
  \item \pp{4} Compute BIC for all four models (3 decimals).
  \item \pp{2} Which model size does each criterion select? Explain the reason for the
    disagreement by comparing the two penalty terms.
  \item \pp{2} Why can the training RSS (or equivalently the training $R^2$) alone never be
    used to select $d$? What does RSS necessarily do as $d$ grows?
\end{enumerate}

\ifdefined\withsolutions
\begin{solutionbox}
\textbf{(a) $C_p$.} Penalty per predictor: $2\hat{\sigma}^2=8$; $C_p=(\operatorname{RSS}_d+8d)/100$:
\begin{center}
\begin{tabular}{lcccc}
\toprule
$d$ & 1 & 2 & 3 & 4\\
\midrule
$\operatorname{RSS}_d+8d$ & 508 & 460 & 452 & 456\\
$C_p$ & 5.080 & 4.600 & \textbf{4.520} & 4.560\\
\bottomrule
\end{tabular}
\end{center}

\medskip
\textbf{(b) BIC.} Penalty per predictor: $\log(100)\,\hat{\sigma}^2=4.605\cdot 4=18.421$;
$\operatorname{BIC}=(\operatorname{RSS}_d+18.421\,d)/100$:
\begin{center}
\begin{tabular}{lcccc}
\toprule
$d$ & 1 & 2 & 3 & 4\\
\midrule
$\operatorname{RSS}_d+18.421\,d$ & 518.421 & 480.841 & 483.262 & 497.683\\
BIC & 5.184 & \textbf{4.808} & 4.833 & 4.977\\
\bottomrule
\end{tabular}
\end{center}

\medskip
\textbf{(c) Selected models.}
$C_p$ is minimised at $d=3$ (4.520); BIC is minimised at $d=2$ (4.808).
Reason: both criteria are RSS plus a penalty for model size, but BIC charges
$\log(n)\hat{\sigma}^2=18.421$ per predictor whereas $C_p$ charges only
$2\hat{\sigma}^2=8$; since $\log(100)=4.605>2$, BIC penalises additional predictors much
more heavily and therefore tends to select \emph{smaller} models. Here the RSS drop from
$d=2$ to $d=3$ (16) exceeds the $C_p$ penalty (8) but not the BIC penalty (18.421).

\medskip
\textbf{(d) Why RSS alone cannot choose $d$.}
For nested model searches, adding a predictor can never increase the training RSS: RSS
decreases \emph{monotonically} in $d$ (and training $R^2$ increases monotonically).
RSS-based selection would therefore always pick the largest model, $d=4$, regardless of
whether the extra predictors carry real signal --- training error is an ever more
optimistic estimate of test error as flexibility grows. Hence one must either penalise
model size ($C_p$, AIC, BIC, adjusted $R^2$) or estimate the test error directly
(validation set, cross-validation).
\end{solutionbox}
\fi

% ============================================================
\problem{6: Ridge regression and the lasso}{15}

Consider a linear model with response $y$ and $p$ standardised predictors
$X_1,\dots,X_p$. Part (e) also draws on the bias--variance ideas of Chapter~2 and the
least-squares fit of Chapter~3.

\begin{enumerate}[label=(\alph*),leftmargin=7mm,itemsep=2.5mm]
  \item \pp{4} Write down the optimisation objectives of \emph{ridge regression} and of
    the \emph{lasso} (RSS plus the respective penalty, with tuning parameter
    $\lambda\ge 0$). Which coefficient is conventionally \emph{not} penalised?
  \item \pp{3} Describe the behaviour of the two estimators in the limits
    $\lambda\to 0$ and $\lambda\to\infty$. Which method from Chapter~3 is recovered as
    $\lambda\to 0$?
  \item \pp{3} Explain --- in words, using the constraint-region geometry for $p=2$ --- why
    the lasso can set coefficients \emph{exactly} to zero while ridge regression
    typically does not.
  \item \pp{2} Why should the predictors be \emph{standardised} before ridge or lasso is
    applied, when this is irrelevant for ordinary least squares?
  \item \pp{3} A ridge regression was fitted for two values of the tuning parameter; the
    table shows the standardised coefficient estimates. A 10-fold cross-validation over a
    grid of $\lambda$ values gave the CV errors on the right.
    \begin{center}
    \begin{tabular}{lrr}
    \toprule
     & $\lambda=0.1$ & $\lambda=100$\\
    \midrule
    $\hat{\beta}_1$ & $0.92$ & $0.31$\\
    $\hat{\beta}_2$ & $0.85$ & $0.29$\\
    $\hat{\beta}_3$ & $-0.41$ & $-0.12$\\
    $\hat{\beta}_4$ & $0.20$ & $0.06$\\
    \bottomrule
    \end{tabular}
    \hspace{12mm}
    \begin{tabular}{lcccc}
    \toprule
    $\lambda$ & 0.1 & 1 & 10 & 100\\
    \midrule
    CV error & 12.8 & 11.9 & 10.6 & 13.4\\
    \bottomrule
    \end{tabular}
    \end{center}
    Which of the two fits shown in the coefficient table has the lower bias and which the
    lower variance, and why? Which $\lambda$ from the CV table should be chosen for the
    final model?
\end{enumerate}

\ifdefined\withsolutions
\begin{solutionbox}
\textbf{(a) Objectives.} Ridge regression minimises
\[
\sum_{i=1}^{n}\Bigl(y_i-\beta_0-\sum_{j=1}^{p}\beta_j x_{ij}\Bigr)^{2}
+\lambda\sum_{j=1}^{p}\beta_j^{2}
=\operatorname{RSS}+\lambda\|\beta\|_2^2 ,
\]
the lasso minimises
\[
\sum_{i=1}^{n}\Bigl(y_i-\beta_0-\sum_{j=1}^{p}\beta_j x_{ij}\Bigr)^{2}
+\lambda\sum_{j=1}^{p}|\beta_j|
=\operatorname{RSS}+\lambda\|\beta\|_1 .
\]
The intercept $\beta_0$ is \emph{not} penalised (shrinking the overall level of $y$ would
make no sense).

\medskip
\textbf{(b) Limiting behaviour.}
$\lambda\to 0$: the penalty vanishes and both estimators approach the \emph{ordinary
least-squares} solution of Chapter~3.
$\lambda\to\infty$: the penalty dominates and all penalised coefficients are forced to
zero --- the null model containing only the intercept. Ridge shrinks the coefficients
\emph{towards} zero (exactly zero only in the limit), whereas the lasso sets more and more
coefficients \emph{exactly} to zero as $\lambda$ grows. In between, $\lambda$ traces out a
whole path of fits whose flexibility decreases as $\lambda$ increases.

\medskip
\textbf{(c) Geometry of the penalties ($p=2$).}
The penalised problems are equivalent to minimising the RSS subject to a budget
constraint: $|\beta_1|+|\beta_2|\le s$ (lasso) or $\beta_1^2+\beta_2^2\le s$ (ridge).
The lasso region is a \emph{diamond} with corners on the coordinate axes; the ridge region
is a \emph{disk} with no corners. The RSS is minimised over the constraint region where
its elliptical contours around the OLS solution first touch the region. Ellipses often hit
the diamond exactly at a \emph{corner} --- and at a corner one coefficient is exactly zero
(sparsity, variable selection). The circle has no corners, so the touching point generically
has both coordinates non-zero: ridge shrinks but keeps all $p$ predictors.

\medskip
\textbf{(d) Standardisation.}
OLS is scale-equivariant: rescaling $X_j$ (e.g.\ EUR $\to$ thousand EUR) simply rescales
$\hat{\beta}_j$, leaving fits unchanged. The penalties, however, treat all coefficients
symmetrically: $\lambda\sum\beta_j^2$ resp.\ $\lambda\sum|\beta_j|$ depend on the
\emph{magnitudes} of the coefficients and hence on the units of the predictors --- a
predictor measured in small units has a large coefficient and would be penalised
(shrunk) more strongly for purely cosmetic reasons. Standardising each predictor to
standard deviation one makes the penalty act comparably on all predictors and the solution
independent of the original units.

\medskip
\textbf{(e) Bias--variance and choice of $\lambda$.}
$\lambda=0.1$: almost no shrinkage --- the fit is close to OLS, i.e.\ more flexible:
\emph{lower bias, higher variance}. $\lambda=100$: strong shrinkage (coefficients roughly
one third of their former size, e.g.\ $0.92\to 0.31$): a more rigid fit with
\emph{higher bias but lower variance} --- the bias--variance trade-off of Chapter~2
governed by a single knob $\lambda$.
From the CV table, the estimated test error is minimised at $\lambda=10$
(CV error $10.6<11.9<12.8<13.4$), so \emph{$\lambda=10$} should be chosen --- an
intermediate amount of shrinkage; neither of the two fits shown in the coefficient table
is optimal.
\end{solutionbox}
\fi

\vspace{6mm}
\begin{center}
\rule{0.35\textwidth}{0.4pt}\\[1mm]
\textit{End of the examination --- good luck!}
\end{center}

\end{document}
